% Seminarium: rozdział 2 pracy
% J. Alman, V. Vassilevska Williams, "Truly Subquadratic 3SUM and Truly Subcubic
% APSP via Triangles in Sparse Lopsided Graphs" (arXiv:2610.06783, 2026).
% Budowanie: ./build.sh  (wersja z animacjami i wersja handout).
\documentclass[aspectratio=169,10pt,t]{beamer}

\usetheme{Madrid}
\usecolortheme{default}
\usefonttheme[onlymath]{serif}
\setbeamertemplate{navigation symbols}{}
\setbeamertemplate{theorems}[numbered]
\setbeamercovered{invisible}

\usepackage[polish]{babel}
\usepackage[T1]{fontenc}
\usepackage{lmodern}
\usepackage{amsmath,amssymb,mathtools}
\usepackage{booktabs}
\usepackage{array}
\usepackage{tikz}
\usetikzlibrary{arrows.meta,positioning,calc,decorations.pathreplacing,fit,
  backgrounds,shapes.misc,babel}
% overlaye w TikZ: [visible on=<2->]
\tikzset{
  invisible/.style={opacity=0,text opacity=0},
  visible on/.style={alt={#1{}{invisible}}},
  alt/.code args={<#1>#2#3}{\alt<#1>{\pgfkeysalso{#2}}{\pgfkeysalso{#3}}}}

% ------------------------------------------------------------------ kolory
\colorlet{zew}{blue!65!black}        % zmienne zewnętrzne (iloczyn zewnętrzny)
\colorlet{wew}{orange!90!black}      % zmienne wewnętrzne (iloczyn skalarny)
\colorlet{chciane}{red!80!black}     % chciane wpisy / odwiedzane liście
\colorlet{szary}{black!35}
\colorlet{tlozew}{blue!12}
\colorlet{tlowew}{orange!20}
\newcommand{\Zew}[1]{\textcolor{zew}{#1}}
\newcommand{\Wew}[1]{\textcolor{wew}{#1}}
\newcommand{\Ch}[1]{\textcolor{chciane}{#1}}

% ------------------------------------------------------------------ twierdzenia z numeracją jak w pracy
\newcommand{\numer}{}
\theoremstyle{plain}
\newtheorem*{twX}{Twierdzenie~\numer}
\newtheorem*{lemX}{Lemat~\numer}

\newenvironment{twierdzenie}[1]{\renewcommand{\numer}{#1}\begin{twX}}{\end{twX}}
\newenvironment{lemat}[1]{\renewcommand{\numer}{#1}\begin{lemX}}{\end{lemX}}
\newenvironment{uwagaP}[1]{\renewcommand{\numer}{#1}\begin{uwX}}{\end{uwX}}
\theoremstyle{definition}
\newtheorem*{uwX}{Uwaga~\numer}
\newtheorem*{definicja}{Definicja}
\newtheorem*{obserwacja}{Obserwacja}
\newtheorem*{wniosek}{Wniosek}
\theoremstyle{example}
\newtheorem*{przyklad}{Przykład}

% ------------------------------------------------------------------ notacja z pracy
\newcommand{\Z}{\mathbb{Z}}
\newcommand{\Full}{\mathsf{Full}}
\newcommand{\Pruned}{\mathsf{Pruned}}
\newcommand{\Mult}{\mathsf{Mult}}
\newcommand{\Leaves}{\mathsf{Leaves}}
\newcommand{\ph}{\hat p}
\newcommand{\qh}{\hat q}
\newcommand{\pusty}{\epsilon}
\newcommand{\idea}[1]{%
  \begin{beamercolorbox}[rounded=true,shadow=true,sep=6pt]{block body alerted}%
  \textbf{\textcolor{chciane}{Kluczowa idea.}} #1\end{beamercolorbox}}

\AtBeginSection[]{%
  \begin{frame}{Plan}
    \tableofcontents[currentsection,hideallsubsections]
  \end{frame}}

\title[Rozdział 2: wybrane wpisy cienkiego iloczynu]{Szybkie liczenie wybranych wpisów\\ cienkiego iloczynu macierzy}
\subtitle{Rozdział 2 pracy J.~Almana i V.~Vassilevskiej Williams (2026)}
\author{[Imię Nazwisko]}
\date{\today}

\begin{document}

\begin{frame}[c]
  \titlepage
  \footnotesize\centering
  J.~Alman, V.~Vassilevska Williams, \emph{Truly Subquadratic 3SUM and Truly Subcubic APSP\\
  via Triangles in Sparse Lopsided Graphs}, arXiv:2610.06783, 2026.
\end{frame}

\begin{frame}{Plan}
  \tableofcontents[hideallsubsections]
\end{frame}

% ======================================================================
\section{Twierdzenie 5: problem i wynik}
% ======================================================================

\begin{frame}{Model obliczeń i problem}
  \begin{columns}[T]
    \begin{column}{0.53\textwidth}
      \begin{itemize}
        \item \textbf{Model:} word RAM ze słowami długości $O(\log N)$ bitów.
          Liczymy \emph{operacje} na liczbach całkowitych o $O(\log N)$ bitach.
        \item \textbf{Dane:} $X\in\Z^{N\times D}$, $Y\in\Z^{D\times N}$, gdzie $D\ll N$.
          Iloczyn $XY$ nazywamy \emph{cienkim}, bo macierze wejściowe są bardzo prostokątne.
        \item Zbiór $W$ pozycji $(I,J)$ macierzy $N\times N$, gdzie $I,J\in\{1,\dots,N\}$.
          Pozycje z $W$ i wpisy $XY$ na nich nazywamy \Ch{\textbf{chcianymi}}.
        \item \textbf{Cel:} policzyć tylko
          \[(XY)[I,J]=\sum_{k=1}^{D}X[I,k]\,Y[k,J]\quad\text{dla }(I,J)\in W.\]
      \end{itemize}
    \end{column}
    \begin{column}{0.45\textwidth}
      \centering
      \begin{tikzpicture}[scale=0.62]
        % X: N x D
        \draw[fill=tlozew] (0,0) rectangle (0.8,5);
        \node at (0.4,-0.35) {\small $D$};
        \node[left] at (0,2.5) {\small $N$};
        \node at (0.4,5.4) {$X$};
        % Y: D x N
        \draw[fill=tlozew] (1.4,4.2) rectangle (6.4,5);
        \node at (3.9,5.4) {$Y$};
        \node[right] at (6.4,4.6) {\small $D$};
        % XY
        \draw (1.4,0) rectangle (6.4,3.8);
        \node at (3.9,-0.35) {\small $XY$: $N\times N$};
        \foreach \x/\y in {1.9/3.2,2.7/1.1,3.3/2.6,4.6/0.6,5.8/3.4,5.2/2.0,2.2/0.5,4.0/3.0,6.0/1.3,3.6/1.6}
          \fill[chciane] (\x,\y) circle (2.2pt);
        \only<2->{
          \fill[chciane!25] (0,2.45) rectangle (0.8,2.75);
          \fill[chciane!25] (4.45,4.2) rectangle (4.75,5);
          \draw[chciane,thick] (0,2.45) rectangle (0.8,2.75);
          \draw[chciane,thick] (4.45,4.2) rectangle (4.75,5);
          \draw[chciane,dashed] (0.8,2.6) -- (4.6,2.6);
          \draw[chciane,dashed] (4.6,4.2) -- (4.6,2.6);
          \fill[chciane] (4.6,2.6) circle (3.2pt);
          \node[chciane,font=\scriptsize,align=center] at (3.2,-1.3)
            {chciany wpis = iloczyn skalarny\\ wiersza $X$ i kolumny $Y$\\ ($D$ mnożeń)};
        }
      \end{tikzpicture}
    \end{column}
  \end{columns}
\end{frame}

\begin{frame}{Główne twierdzenie rozdziału}
  \begin{twierdzenie}{5}
    Niech $D\ge4$ będzie potęgą czwórki i niech $N\ge D^{18}$. Dla danych macierzy
    $X\in\Z^{N\times D}$ i $Y\in\Z^{D\times N}$, których wpisy są liczbami całkowitymi
    o wartości bezwzględnej co najwyżej $N^{O(1)}$, oraz zbioru $W$ co najwyżej
    $N^2/\sqrt D$ pozycji macierzy $N\times N$, wpisy $(XY)[I,J]$ dla $(I,J)\in W$
    można obliczyć deterministycznie w czasie
    \[
      O\!\left(\frac{N^2\log^2 D}{D^{1/18}}\right).
    \]
  \end{twierdzenie}
  \pause
  \begin{itemize}
    \item $X$, $Y$ i $XY$ mogą być \textbf{dowolnymi gęstymi} macierzami.
    \item O zbiorze $W$ zakładamy tylko tyle, że $|W|\le N^2/\sqrt D$. Nie zakładamy żadnej struktury.
    \item Stałe (np. wykładnik $18$) dobrano tak, żeby prezentacja była prosta.
  \end{itemize}
\end{frame}

\begin{frame}{Z czym się porównujemy?}
  \centering
  \renewcommand{\arraystretch}{1.5}
  \begin{tabular}{lll}
    \toprule
    \textbf{Metoda} & \textbf{Liczba operacji} & \textbf{Uzasadnienie}\\
    \midrule
    Policzyć (wypisać) całe $XY$ & $\ge N^2$ & $XY$ ma $N^2$ wpisów\\
    \uncover<2->{Każdy chciany wpis osobno} & \uncover<2->{$|W|\cdot D\le \dfrac{N^2}{\sqrt D}\cdot D=N^2\sqrt D$} &
      \uncover<2->{$D$ mnożeń na wpis}\\
    \uncover<3->{\textbf{Twierdzenie 5}} & \uncover<3->{$O\!\left(N^2\log^2 D\cdot D^{-1/18}\right)$} & \uncover<3->{ten rozdział}\\
    \bottomrule
  \end{tabular}
  \vspace{1em}
  \uncover<4->{
  \begin{itemize}
    \item Wynik jest \textbf{wielomianowo} (w $D$) mniejszy niż \emph{wypisanie} $XY$ i niż
      $N^2/\sqrt D$ osobnych iloczynów skalarnych.
    \item Przykład: przy największym dozwolonym $D=N^{1/18}$ mamy $D^{1/18}=N^{1/324}$, czyli
      czas $O(N^{2-1/324}\log^2 N)$.
  \end{itemize}}
\end{frame}

\begin{frame}{Mapa rozdziału}
  \centering
  \begin{tikzpicture}[
      box/.style={draw,rounded corners,fill=blue!8,align=center,text width=2.7cm,
        minimum height=2.3cm,font=\small},
      arr/.style={-{Stealth[length=3mm]},thick}]
    \node[box] (s) {\textbf{2.1 Strassen}\\ rekurencja,\\ kodowanie\,/\\ dekodowanie};
    \node[box,right=0.6cm of s] (h) {\textbf{2.2 Schönhage}\\ iloczyn zewnętrzny\\ $+$ skalarny\\ $13\to10$ mnożeń};
    \node[box,right=0.6cm of h] (f) {\textbf{2.3 $\Full$}\\ $L$ poziomów:\\ wiele iloczynów\\ naraz, kafelki};
    \node[box,right=0.6cm of f,fill=red!10] (p) {\textbf{2.4 $\Pruned$}\\ \emph{nowe}: tylko\\ liście potrzebne\\ dla $W$};
    \draw[arr] (s)--(h); \draw[arr] (h)--(f); \draw[arr] (f)--(p);
  \end{tikzpicture}
  \vspace{0.6em}
  \begin{itemize}
    \item<2-> Mnożenia odbywają się w \textbf{liściach} drzewa rekursji. Dodawania na drodze w dół
      (\emph{kodowanie}) i w górę (\emph{dekodowanie}) są stosunkowo tanie.
    \item<3-> \textbf{Nowość 1:} kodowanie zależy \emph{tylko od $X$} albo \emph{tylko od $Y$},
      więc liczymy je raz i współdzielimy między wieloma wywołaniami.
    \item<4-> \textbf{Nowość 2:} przycinamy drzewo i odwiedzamy tylko liście potrzebne chcianym wpisom.
      Każdy wpis ma \emph{prywatny} liść, a pozostałe jego liście są współdzielone.
      Wychodzi $\approx N^2/D^{1/18}$ liści po $O(\log^2 D)$ operacji każdy.
  \end{itemize}
\end{frame}

% ======================================================================
\section{2.1 Rekurencyjny algorytm Strassena}
% ======================================================================

\begin{frame}{Tożsamość Strassena: macierze $2\times2$ przez 7 mnożeń}
  Niech $A=(a_{ij})$ i $B=(b_{ij})$ będą macierzami $2\times2$. Tworzymy siedem iloczynów:
  \begin{equation*}
    \begin{aligned}
      S_1&=(a_{11}+a_{22})(b_{11}+b_{22}), & S_5&=(a_{11}+a_{12})\,b_{22},\\
      S_2&=(a_{21}+a_{22})\,b_{11},        & S_6&=(a_{21}-a_{11})(b_{11}+b_{12}),\\
      S_3&=a_{11}(b_{12}-b_{22}),          & S_7&=(a_{12}-a_{22})(b_{21}+b_{22}),\\
      S_4&=a_{22}(b_{21}-b_{11}),
    \end{aligned}\tag{1}
  \end{equation*}
  \pause
  i wtedy
  \[
    AB=\begin{pmatrix} S_1+S_4-S_5+S_7 & S_3+S_5\\ S_2+S_4 & S_1-S_2+S_3+S_6\end{pmatrix}.
  \]
  \pause
  \begin{itemize}
    \item Każde $S_r$ to \Zew{(kombinacja liniowa wpisów $A$)} $\times$ \Wew{(kombinacja liniowa wpisów $B$)}.
    \item Każdy wpis $AB$ to kombinacja liniowa $S_1,\dots,S_7$.
    \item \textbf{7 mnożeń zamiast 8.}
  \end{itemize}
\end{frame}

\begin{frame}{Rekurencja na blokach}
  \begin{columns}[T]
    \begin{column}{0.42\textwidth}
      \small
      \begin{itemize}
        \item $A,B$ rozmiaru $N\times N$, gdzie $N=2^L$.
        \item Traktujemy je jak macierze $2\times2$ \emph{bloków} rozmiaru $\frac N2\times\frac N2$.
        \item<2-> Tożsamość nie korzysta z przemienności mnożenia, więc działa dla bloków.
        \item<3-> Tworzymy 7 par kombinacji bloków, każdą parę mnożymy \emph{rekurencyjnie},
          a z 7 wyników składamy 4 bloki $AB$.
        \item<4-> Drzewo ma $L$ poziomów i $7^L$ liści, po jednym na każdy wybór jednego z 7
          iloczynów na każdym poziomie.
        \item<5-> Mnożenia są \textbf{w liściach}:
          \[7^L=N^{\log_27}<N^{2.81}.\]
      \end{itemize}
    \end{column}
    \begin{column}{0.56\textwidth}
      \centering
      \begin{tikzpicture}[
          v/.style={draw,rounded corners=2pt,fill=blue!8,font=\scriptsize,inner sep=2pt,minimum width=0.62cm},
          scale=0.7,transform shape]
        \node[v] (r) at (3.5,0) {$N\times N$};
        \foreach \i in {1,...,7} {
          \node<3->[v] (c\i) at ({(\i-1)*1.15},-1.5) {$\frac N2{\times}\frac N2$};
          \draw<3-> (r.south) -- (c\i.north);
        }
        \foreach \i in {1,...,7} {
          \foreach \j in {1,...,7} {
            \draw<4->[szary] (c\i.south) -- ({(\i-1)*1.15+(\j-4)*0.14},-2.6);
          }
          \node<4->[font=\scriptsize] at ({(\i-1)*1.15},-2.9) {$\vdots$};
        }
        \node<4->[font=\small] at (3.5,-3.6) {$7^L$ liści: iloczyny dwóch liczb};
        \draw<3->[-{Stealth},thick,zew] (-0.9,-0.1) -- node[left,align=right,font=\scriptsize]
          {w dół:\\ kodowanie} (-0.9,-2.8);
        \draw<3->[-{Stealth},thick,wew] (7.9,-2.8) -- node[right,align=left,font=\scriptsize]
          {w górę:\\ dekodowanie} (7.9,-0.1);
      \end{tikzpicture}

      \uncover<5->{\footnotesize Równoważnie: $T(N)=7\,T(N/2)+O(N^2)$.}
    \end{column}
  \end{columns}
\end{frame}

\begin{frame}{Koszt dodawań i odejmowań}
  \begin{itemize}
    \item Na głębokości $k$ drzewa jest $7^k$ wierzchołków wewnętrznych.
    \item W każdym z nich, w dół i w górę, liczymy kombinacje liniowe macierzy
      $2^{L-k}\times2^{L-k}$, czyli $O(4^{L-k})$ operacji.
  \end{itemize}
  \pause
  \[
    \sum_{k=0}^{L-1}7^k4^{L-k}
    \;=\;4^L\sum_{k=0}^{L-1}\Big(\tfrac74\Big)^k
    \;=\;4^L\cdot\frac{(7/4)^L-1}{7/4-1}
    \;<\;\tfrac43\cdot7^L
    \;=\;O(7^L).
  \]
  \pause
  Szereg geometryczny jest zdominowany przez ostatni wyraz, więc dodawania zwiększają czas tylko o stały czynnik.
  \pause
  \alert{Uwaga:} w nowym algorytmie (2.4) ten argument \emph{przestanie działać}. Dodawania trzeba
  będzie liczyć staranniej.
\end{frame}

\begin{frame}{Wnioski: trzy fazy}
  \centering
  \begin{tikzpicture}[
      b/.style={draw,rounded corners,align=center,font=\small,minimum height=1cm,text width=2.5cm},
      arr/.style={-{Stealth[length=2.5mm]},thick}]
    \node (A) at (0,0.8) {$A$};
    \node (B) at (0,-0.8) {$B$};
    \node[b,fill=tlozew] (kA) at (2.6,0.8) {kodowanie $A$\\ \scriptsize liniowe, tylko z $A$};
    \node[b,fill=tlozew] (kB) at (2.6,-0.8) {kodowanie $B$\\ \scriptsize liniowe, tylko z $B$};
    \node[b,fill=red!10,minimum height=2.6cm] (m) at (6.1,0) {mnożenie\\ $7^L$ iloczynów\\ w liściach};
    \node[b,fill=tlowew,minimum height=2.6cm] (d) at (9.4,0) {dekodowanie\\ \scriptsize liniowe};
    \node (AB) at (11.6,0) {$AB$};
    \draw[arr] (A)--(kA); \draw[arr] (B)--(kB);
    \draw[arr] (kA)--(m.west|-kA); \draw[arr] (kB)--(m.west|-kB);
    \draw[arr] (m)--(d); \draw[arr] (d)--(AB);
  \end{tikzpicture}
  \vspace{0.5em}
  \begin{itemize}
    \item<2-> Kodowanie i dekodowanie są \textbf{liniowe}.
    \item<3-> Kodowanie $A$ \textbf{nie zależy od $B$} (i odwrotnie). Wykorzystamy to w 2.4.1.
    \item<4-> Tutaj kodowanie i dekodowanie są tańsze niż mnożenia, więc można je zignorować.
      \alert{W pełnym algorytmie (2.4) już nie!}
  \end{itemize}
  \uncover<5->{\idea{Mała tożsamość stosowana rekurencyjnie: mnożenia w liściach, a ich liczba decyduje o koszcie.}}
\end{frame}

\begin{frame}{Notacja: wielomiany trójliniowe}
  \begin{itemize}
    \item Wprowadzamy \emph{zmienne wyjściowe} $c_{11},c_{12},c_{21},c_{22}$, po jednej na każdy wpis $AB$.
    \item Każde $S_r$ mnożymy przez sumę ze znakami tych $c_{ik}$, w których wzorach $S_r$ występuje.
      Na przykład $S_1$ jest we wzorach na $(AB)_{11}$ i $(AB)_{22}$, więc dostaje $c_{11}+c_{22}$.
      $S_5$ jest z minusem w $(AB)_{11}$ i z plusem w $(AB)_{12}$, więc dostaje $c_{12}-c_{11}$.
  \end{itemize}
  \pause
  \[
    S_1(c_{11}+c_{22})+S_2(c_{21}-c_{22})+S_3(c_{12}+c_{22})+S_4(c_{11}+c_{21})
    +S_5(c_{12}-c_{11})+S_6c_{22}+S_7c_{11}
    =\sum_{i,j,k=1}^{2}a_{ij}b_{jk}c_{ik}.
  \]
  \pause
  \begin{itemize}
    \item Współczynnik przy $c_{ik}$ \textbf{po prawej} to szukane $(AB)_{ik}=a_{i1}b_{1k}+a_{i2}b_{2k}$.
    \item \textbf{Po lewej} to kombinacja, która go liczy. Dla $c_{11}$ jest to $S_1+S_4-S_5+S_7$.
    \item<4-> Każdy składnik ma postać
      $\Zew{\varphi(a)}\cdot\Wew{\psi(b)}\cdot\chi(c)$, czyli (forma liniowa w $a$) $\times$ (forma w $b$) $\times$ (forma w $c$).
      Jedna tożsamość zapisuje \emph{kodowanie}, \emph{mnożenia} i \emph{dekodowanie} naraz.
    \item<4-> Wielomian jest \emph{trójliniowy}: liniowy osobno w $a$, w $b$ i w $c$.
      W literaturze nazywa się go tensorem rzędu~3.
  \end{itemize}
\end{frame}

% ======================================================================
\section{2.2 Tożsamość Schönhagego}
% ======================================================================

\begin{frame}{Dwa iloczyny na rozłącznych zmiennych}
  \begin{columns}[T]
    \begin{column}{0.5\textwidth}
      \textbf{\Zew{Iloczyn zewnętrzny}} wektorów $(x_1,x_2,x_3)$ i $(y_1,y_2,y_3)$:
      \[\text{dziewięć liczb } x_iy_j,\quad 1\le i,j\le3.\]
      Jest to iloczyn macierzy $3\times1$ przez $1\times3$.
      \begin{center}
        \begin{tikzpicture}[scale=0.32]
          \begin{scope}[shift={(0,0)}]\draw[fill=tlozew] (0,0) rectangle (1,3); \draw (0,0) grid (1,3);\end{scope} \node at (0.5,-0.5) {\scriptsize $3{\times}1$};
          \node at (1.6,1.5) {$\cdot$};
          \begin{scope}[shift={(2.2,1)}]\draw[fill=tlozew] (0,0) rectangle (3,1); \draw (0,0) grid (3,1);\end{scope} \node at (3.7,0.5) {\scriptsize $1{\times}3$};
          \node at (5.8,1.5) {$=$};
          \begin{scope}[shift={(6.4,0)}]\draw[fill=tlozew] (0,0) rectangle (3,3); \draw (0,0) grid (3,3);\end{scope} \node at (7.9,-0.5) {\scriptsize $3{\times}3$};
        \end{tikzpicture}
      \end{center}
    \end{column}
    \begin{column}{0.5\textwidth}
      \textbf{\Wew{Iloczyn skalarny}} wektorów $(p_{11},p_{12},p_{21},p_{22})$ i $(q_{11},q_{12},q_{21},q_{22})$:
      \[\textstyle\sum_{i,j=1}^{2}p_{ij}q_{ij}.\]
      Jest to iloczyn macierzy $1\times4$ przez $4\times1$.
      \begin{center}
        \begin{tikzpicture}[scale=0.32]
          \begin{scope}[shift={(0,1)}]\draw[fill=tlowew] (0,0) rectangle (4,1); \draw (0,0) grid (4,1);\end{scope} \node at (2,0.5) {\scriptsize $1{\times}4$};
          \node at (4.6,1.5) {$\cdot$};
          \begin{scope}[shift={(5.2,-0.5)}]\draw[fill=tlowew] (0,0) rectangle (1,4); \draw (0,0) grid (1,4);\end{scope} \node at (5.7,-1) {\scriptsize $4{\times}1$};
          \node at (6.8,1.5) {$=$};
          \draw[fill=tlowew] (7.4,1) rectangle (8.4,2); \node at (7.9,0.5) {\scriptsize $1{\times}1$};
        \end{tikzpicture}
      \end{center}
      {\scriptsize Indeksy $\{1,2\}^2$ zamiast $\{1,2,3,4\}$ ułatwią zapis $\ph$, $\qh$.}
    \end{column}
  \end{columns}
  \pause
  \begin{itemize}\small
    \item Osobno: $9+4=\mathbf{13}$ mnożeń.
    \item<3-> \textbf{Schönhage [Sch81]:} wystarczy \alert{\textbf{10}} mnożeń, jeśli dopuścimy pewne
      wyrazy błędu (okażą się nieszkodliwe).
  \end{itemize}
\end{frame}

\begin{frame}{Zapis trójliniowy: wielomian $G$}
  \begin{itemize}
    \item \textbf{7 zmiennych lewych:} $\Zew{x_1,x_2,x_3}$, $\Wew{p_{11},p_{12},p_{21},p_{22}}$.
    \item \textbf{7 zmiennych prawych:} $\Zew{y_1,y_2,y_3}$, $\Wew{q_{11},q_{12},q_{21},q_{22}}$.
    \item \textbf{10 zmiennych wyjściowych:} $\Zew{z_{ij}}$ ($1\le i,j\le3$) i $\Wew{z_0}$, po jednej na każdą z dziesięciu liczb do policzenia.
  \end{itemize}
  \pause
  \[
    G:=\sum_{i,j=1}^{3}\Zew{x_iy_j\,z_{ij}}\;+\;\Big(\sum_{i,j=1}^{2}\Wew{p_{ij}q_{ij}}\Big)\Wew{z_0}.
  \]
  \pause
  \begin{itemize}
    \item Współczynnik przy $z_{ij}$ to $x_iy_j$ (iloczyn zewnętrzny).
    \item Współczynnik przy $z_0$ to $\sum p_{ij}q_{ij}$ (iloczyn skalarny).
    \item Chcemy zapisać $G$ (z dokładnością do błędu) jako sumę \textbf{dziesięciu} wyrazów
      postaci (forma liniowa w zmiennych lewych)$\times$(forma w prawych)$\times$(forma w wyjściowych).
  \end{itemize}
\end{frame}

\begin{frame}{Macierze $\ph$ i $\qh$}
  Rozszerzamy macierze $2\times2$ zmiennych $p$ i $q$ do $3\times3$ tak, żeby
  \textbf{każda kolumna $\ph$} i \textbf{każdy wiersz $\qh$} sumowały się do zera:
  \[
    \ph=\begin{pmatrix} p_{11}&p_{12}&0\\ p_{21}&p_{22}&0\\ -p_{11}-p_{21}&-p_{12}-p_{22}&0\end{pmatrix},
    \qquad
    \qh=\begin{pmatrix} q_{11}&q_{12}&-q_{11}-q_{12}\\ q_{21}&q_{22}&-q_{21}-q_{22}\\ 0&0&0\end{pmatrix}.
  \]
  \pause
  \begin{itemize}
    \item Elementy $\ph_{ij}$, $\qh_{ij}$ to \emph{formy liniowe}, a nie nowe zmienne.
    \item $\displaystyle\sum_{i,j=1}^{3}\ph_{ij}\qh_{ij}=\sum_{i,j=1}^{2}p_{ij}q_{ij}$, bo trzecia
      kolumna $\ph$ i trzeci wiersz $\qh$ są zerowe, a dla $i,j\le2$ mamy $\ph_{ij}=p_{ij}$, $\qh_{ij}=q_{ij}$.
  \end{itemize}
\end{frame}

\begin{frame}{Dziesięć składników i ich formy liniowe}
  Składniki to $P_{ij}$ (dla $1\le i,j\le3$, odpowiadające wpisom $\ph$ i $\qh$) oraz $P_0$.
  Każdy \textbf{składnik} $\lambda$ jest iloczynem trzech form liniowych:
  $\varphi_\lambda$ w zmiennych lewych, $\psi_\lambda$ w prawych, $\chi_\lambda$ w wyjściowych.
  \vspace{0.4em}
  \begin{center}
    \renewcommand{\arraystretch}{1.4}
    \begin{tabular}{c|c|c|c}
      składnik $\lambda$ & $\varphi_\lambda$ (lewe) & $\psi_\lambda$ (prawe) & $\chi_\lambda$ (wyjściowe)\\\hline
      $P_{ij}$, $1\le i,j\le3$ & $x_i+\ph_{ij}$ & $y_j+\qh_{ij}$ & $z_{ij}+z_0$\\
      $P_0$ & $-(x_1+x_2+x_3)$ & $y_1+y_2+y_3$ & $z_0$
    \end{tabular}
  \end{center}
  \pause
  \begin{itemize}
    \item 10 składników to \textbf{10 mnożeń} $\varphi_\lambda\cdot\psi_\lambda$.
    \item Każde $\varphi_\lambda$ i każde $\psi_\lambda$ ma \textbf{co najwyżej 3} jednomiany o współczynnikach $\pm1$.
  \end{itemize}
\end{frame}

\begin{frame}{Lemat 6: tożsamość Schönhagego}
  \begin{lemat}{6}
    Sumując po dziesięciu składnikach,
    \[
      \sum_\lambda\varphi_\lambda\psi_\lambda\chi_\lambda=G+E,
      \qquad\text{gdzie}\qquad
      E:=\sum_{i,j=1}^{3}\big(x_i\qh_{ij}+\ph_{ij}y_j+\ph_{ij}\qh_{ij}\big)z_{ij}.
    \]
  \end{lemat}
  \pause
  \begin{itemize}
    \item Lewa strona: 10 wyrazów, czyli \textbf{10 mnożeń} $\varphi_\lambda\cdot\psi_\lambda$.
    \item Współczynnik przy $z_0$ po lewej to \textbf{dokładnie} iloczyn skalarny $\sum p_{ij}q_{ij}$.
    \item Współczynnik przy $z_{ij}$ to $x_iy_j$ \textbf{plus błąd} $x_i\qh_{ij}+\ph_{ij}y_j+\ph_{ij}\qh_{ij}$.
    \item<3-> Lemat 6 to kluczowa tożsamość: liczy chciane $G$ plus wielomian błędu $E$.
  \end{itemize}
\end{frame}

\begin{frame}{Dowód Lematu 6}
  \begin{proof}
    Porównujemy współczynniki przy każdej zmiennej wyjściowej po lewej stronie.
    \begin{itemize}
      \item<2-> \textbf{Przy $z_{ij}$:} $z_{ij}$ występuje tylko w $\chi_{P_{ij}}$, więc współczynnik wynosi
        \[\varphi_{P_{ij}}\psi_{P_{ij}}=(x_i+\ph_{ij})(y_j+\qh_{ij})
          =x_iy_j+\big(x_i\qh_{ij}+\ph_{ij}y_j+\ph_{ij}\qh_{ij}\big).\]
      \item<3-> \textbf{Przy $z_0$:} $z_0$ występuje we wszystkich $\chi_\lambda$, więc współczynnik wynosi
        \[
          \sum_{i,j=1}^{3}(x_i+\ph_{ij})(y_j+\qh_{ij})-(x_1+x_2+x_3)(y_1+y_2+y_3)
          =\underbrace{\sum_{i}x_i\sum_{j}\qh_{ij}}_{=0\ (\text{wiersze }\qh)}
           +\underbrace{\sum_{j}y_j\sum_{i}\ph_{ij}}_{=0\ (\text{kolumny }\ph)}
           +\sum_{i,j=1}^{3}\ph_{ij}\qh_{ij}.
        \]
        Iloczyny $x_iy_j$ się skracają, a $\sum_{i,j=1}^{3}\ph_{ij}\qh_{ij}=\sum_{i,j=1}^{2}p_{ij}q_{ij}$. \qedhere
    \end{itemize}
  \end{proof}
\end{frame}

\begin{frame}{Obserwacja 1: zmienne wewnętrzne i zewnętrzne}
  \begin{definicja}
    Zmienne iloczynu skalarnego, czyli $\Wew{p_{ij},q_{ij},z_0}$, nazywamy \Wew{\textbf{wewnętrznymi}}.
    Zmienne iloczynu zewnętrznego, czyli $\Zew{x_i,y_j,z_{ij}}$, nazywamy \Zew{\textbf{zewnętrznymi}}.
  \end{definicja}
  \pause
  \begin{obserwacja}
    \begin{itemize}
      \item Każdy jednomian $G$ jest \emph{czysto} wewnętrzny ($\Wew{p_{ij}q_{ij}z_0}$) albo
        \emph{czysto} zewnętrzny ($\Zew{x_iy_jz_{ij}}$).
      \item Każdy jednomian $E$ ma \Zew{zewnętrzną zmienną wyjściową $z_{ij}$} razem
        z \Wew{wewnętrzną zmienną lewą lub prawą}:
        $\Zew{x_i}\Wew{\qh_{ij}}\Zew{z_{ij}}$, \ $\Wew{\ph_{ij}}\Zew{y_jz_{ij}}$, \ $\Wew{\ph_{ij}\qh_{ij}}\Zew{z_{ij}}$.
    \end{itemize}
  \end{obserwacja}
  To pozwoli pokazać, że $E$ jest nieszkodliwe (Lemat 9).
\end{frame}

\begin{frame}{Obserwacja 2: kto wnosi wkład do którego wyjścia}
  \begin{columns}[T]
    \begin{column}{0.5\textwidth}
      \small
      \begin{definicja}
        Składnik $\lambda$ \textbf{wnosi wkład} do zmiennej wyjściowej $z$, jeśli $z$ występuje w $\chi_\lambda$.
      \end{definicja}
      \begin{itemize}
        \item<2-> \textbf{Każdy} składnik wnosi wkład do $z_0$.
        \item<2-> Do $z_{ij}$ wnosi wkład \textbf{tylko} $P_{ij}$.
        \item<3-> Formy $\chi_\lambda$ są sumami zmiennych bez współczynników, więc współczynnik przy $z$ po lewej stronie to
          \[\sum_{\lambda\text{ wnosi wkład do }z}\varphi_\lambda\psi_\lambda.\]
        \item<4-> To, że do $z_{ij}$ wnosi wkład tylko jeden składnik, pozwoli
          \alert{pominąć wiele mnożeń}, gdy chcemy tylko niektórych wyjść.
      \end{itemize}
    \end{column}
    \begin{column}{0.48\textwidth}
      \centering
      \begin{tikzpicture}[
          t/.style={draw,rounded corners=1pt,fill=red!10,font=\tiny,minimum width=0.5cm,inner sep=1.5pt},
          o/.style={draw,rounded corners=1pt,fill=green!12,font=\tiny,minimum width=0.5cm,inner sep=1.5pt}]
        \foreach \n [count=\k] in {11,12,13,21,22,23,31,32,33} {
          \node[t] (P\n) at ({(\k-1)*0.62},2) {$P_{\n}$};
          \node[o] (z\n) at ({(\k-1)*0.62},0) {$z_{\n}$};
          \draw[thick] (P\n) -- (z\n);
        }
        \node[t] (P0) at (5.9,2) {$P_0$};
        \node[o] (z0) at (5.9,0) {$z_0$};
        \foreach \n in {11,12,13,21,22,23,31,32,33} \draw[szary] (P\n.south) -- (z0.north);
        \draw[thick] (P0) -- (z0);
        \node[font=\scriptsize] at (2.9,2.6) {składniki $\lambda$};
        \node[font=\scriptsize] at (2.9,-0.6) {zmienne wyjściowe};
      \end{tikzpicture}
    \end{column}
  \end{columns}
\end{frame}

% ======================================================================
\section{2.3 Mnożenie prostokątne tożsamością Schönhagego (Coppersmith)}
% ======================================================================

\begin{frame}{Idea: $L$ poziomów rekursji daje $2^L$ iloczynów}
  \begin{itemize}
    \item Tożsamość Schönhagego liczy dwa iloczyny na rozłącznych zmiennych:
      \Zew{zewnętrzny} (kształt $3\times1$ razy $1\times3$) i \Wew{skalarny} ($1\times4$ razy $4\times1$).
    \item Stosowana rekurencyjnie na $L$ poziomach daje $2^L$ iloczynów macierzy, po jednym na każdy wybór
      \Zew{zewnętrzny}/\Wew{skalarny} na każdym poziomie.
    \item<2-> Poziom \Zew{zewnętrzny} mnoży liczbę wierszy i kolumn wyniku przez $3$. Poziom \Wew{skalarny}
      mnoży wymiar wewnętrzny (długość sum) przez $4$.
    \item<3-> Wybór skalarnego na $m$ poziomach i zewnętrznego na pozostałych $L-m$ daje iloczyn macierzy
      \[3^{L-m}\times4^m\quad\text{przez}\quad4^m\times3^{L-m}.\]
      Interesuje nas $m$ \textbf{dużo mniejsze} od $L$, czyli cienkie iloczyny.
  \end{itemize}
  \uncover<4->{\footnotesize
  \begin{przyklad}[$L=2$]
    \centering
    \renewcommand{\arraystretch}{1.2}
    \begin{tabular}{cccc}
      \Zew{zewn.}, \Zew{zewn.} & \Zew{zewn.}, \Wew{skal.} & \Wew{skal.}, \Zew{zewn.} & \Wew{skal.}, \Wew{skal.}\\
      $9\times1$ razy $1\times9$ & $3\times4$ razy $4\times3$ & $3\times4$ razy $4\times3$ & $1\times16$ razy $16\times1$\\
      ($m=0$) & ($m=1$) & ($m=1$) & ($m=2$)
    \end{tabular}
  \end{przyklad}}
\end{frame}

\begin{frame}{Notacja: napisy i tablice}
  \small
  \begin{definicja}
    \begin{itemize}
      \item \textbf{Lewy napis} długości $L$ to ciąg $u=u_1u_2\cdots u_L$ złożony z $L$ zmiennych lewych.
        $u_\ell$ to jego zmienna na \textbf{poziomie} $\ell$.
      \item \textbf{Tablica} $a$ na lewych napisach długości $L$ przypisuje każdemu takiemu $u$ liczbę
        całkowitą $a[u]$. Ma ona $7^L$ wpisów.
      \item Tak samo definiujemy napisy i tablice \textbf{prawe} ($7^L$ wpisów) i \textbf{wyjściowe} ($10^L$ wpisów).
    \end{itemize}
  \end{definicja}
  \pause
  \begin{definicja}
    Dla składnika $\lambda$ i zmiennej lewej $s$ niech $\varphi_\lambda(s)\in\{0,\pm1\}$ oznacza współczynnik
    przy $s$ w $\varphi_\lambda$. Analogicznie $\psi_\lambda(t)$ dla zmiennej prawej $t$.
  \end{definicja}
\end{frame}

\begin{frame}{Slice: odpowiednik bloków macierzy}
  \begin{columns}[T]
    \begin{column}{0.55\textwidth}
      \begin{definicja}
        Dla tablicy $a$ na napisach długości $L\ge1$ i zmiennej lewej $s$ \textbf{slice} $a_s$ to tablica
        na napisach długości $L-1$:
        \[a_s[u']:=a[su'].\]
      \end{definicja}
      \begin{itemize}
        \item $a$ składa się z \textbf{siedmiu} slice'ów, tak jak macierz blokowa składa się z bloków.
        \item Analogicznie slice'y $b_t$ tablicy prawej i $c_z$ tablicy wyjściowej (\textbf{dziesięć}
          slice'ów).
        \item Slice wybieramy według zmiennej na \textbf{poziomie 1}.
      \end{itemize}
    \end{column}
    \begin{column}{0.43\textwidth}
      \centering
      \begin{tikzpicture}[scale=0.36]
        \foreach \s [count=\i] in {x_1,x_2,x_3,p_{11},p_{12},p_{21},p_{22}} {
          \node[font=\tiny,anchor=east] at (0,{7.5-\i}) {$\s$};
          \node[font=\tiny,rotate=90,anchor=west] at ({\i-0.5},7.1) {$\s$};
        }
        \draw[step=1,szary] (0,0) grid (7,7);
        \only<2->{\fill[blue!25] (0,5) rectangle (7,6); \draw[thick,zew] (0,5) rectangle (7,6);
          \node[zew,font=\scriptsize,anchor=west] at (7.1,5.5) {slice $a_{x_2}$};}
        \draw[thick] (0,0) rectangle (7,7);
        \node[font=\scriptsize] at (3.5,-0.6) {$u_2$ (poziom 2)};
        \node[font=\scriptsize,rotate=90] at (-2.1,3.5) {$u_1$ (poziom 1)};
        \node[font=\scriptsize,align=center] at (3.5,-2.1) {tablica $a$ dla $L=2$: $7^2=49$ wpisów};
      \end{tikzpicture}
    \end{column}
  \end{columns}
\end{frame}

\begin{frame}{Algorytm $\Full(a,b)$}
  \begin{block}{$\Full(a,b)$ dla tablic $a,b$ na lewych i prawych napisach długości $L$; zwraca tablicę na wyjściowych napisach długości $L$}\small
    \begin{enumerate}
      \item Jeśli $L=0$, zwróć liczbę $ab$.
      \item \textbf{(kodowanie)} Dla każdego składnika $\lambda$ utwórz
        \[A_\lambda:=\sum_s\varphi_\lambda(s)\,a_s,\qquad B_\lambda:=\sum_t\psi_\lambda(t)\,b_t.\]
      \item \textbf{(rekurencja)} Dla każdego $\lambda$ policz rekurencyjnie $C_\lambda:=\Full(A_\lambda,B_\lambda)$.
      \item \textbf{(dekodowanie)} Zwróć tablicę $c$ o slice'ach
        \[c_{z_{ij}}:=C_{P_{ij}},\qquad c_{z_0}:=\sum_\lambda C_\lambda.\]
    \end{enumerate}
  \end{block}
  \pause
  \begin{itemize}\small
    \item Dla $L=1$ jest to dokładnie jedno zastosowanie tożsamości Schönhagego: krok (2) liczy formy
      $\varphi_\lambda$, $\psi_\lambda$, krok (3) mnoży, a krok (4) sumuje według $\chi_\lambda$.
    \item Krok (4) to $\chi_\lambda$: do $z_{ij}$ wnosi wkład tylko $P_{ij}$, a do $z_0$ wszystkie.
  \end{itemize}
\end{frame}

\begin{frame}{Jeden poziom rekursji}
  \centering
  \begin{tikzpicture}[
      sl/.style={draw,fill=blue!10,font=\scriptsize,minimum width=0.95cm,minimum height=0.5cm,inner sep=1pt},
      en/.style={draw,fill=blue!20,font=\scriptsize,minimum width=0.95cm,minimum height=0.5cm,inner sep=1pt},
      re/.style={draw,fill=orange!20,font=\scriptsize,minimum width=0.95cm,minimum height=0.5cm,inner sep=1pt},
      de/.style={draw,fill=green!15,font=\scriptsize,minimum width=0.95cm,minimum height=0.5cm,inner sep=1pt},
      lab/.style={font=\scriptsize,align=right,anchor=east}]
    \foreach \s/\t [count=\i] in {x1/x_1,x2/x_2,x3/x_3,p11/p_{11},p12/p_{12},p21/p_{21},p22/p_{22}}
      \node[sl] (a\s) at ({1.7+(\i-1)*1.25},3.7) {$a_{\t}$};
    \node[font=\scriptsize] at (5.45,4.25) {tablica $a$ pocięta na 7 slice'ów $a_s$ według zmiennej $s$ na poziomie 1};
    \foreach \l/\t [count=\i] in {P11/P_{11},P12/P_{12},P13/P_{13},P21/P_{21},P22/P_{22},P23/P_{23},P31/P_{31},P32/P_{32},P33/P_{33},P0/P_0} {
      \node[en,visible on=<2->] (A\l) at ({(\i-1)*1.2},2.35) {$A_{\t}$};
      \node[re,visible on=<3->] (C\l) at ({(\i-1)*1.2},1.15) {$C_{\t}$};
      \draw[visible on=<3->] (A\l) -- (C\l);
    }
    \foreach \z/\t [count=\i] in {11/z_{11},12/z_{12},13/z_{13},21/z_{21},22/z_{22},23/z_{23},31/z_{31},32/z_{32},33/z_{33},0/z_0}
      \node[de,visible on=<4->] (c\z) at ({(\i-1)*1.2},0) {$c_{\t}$};
    % kodowanie: + linia ciągła, - przerywana
    \begin{scope}[visible on=<2->]
      \foreach \l/\s in {P11/x1,P11/p11,P12/x1,P12/p12,P13/x1,P21/x2,P21/p21,P22/x2,P22/p22,P23/x2,P31/x3,P32/x3,P33/x3}
        \draw[szary] (a\s.south) -- (A\l.north);
      \foreach \l/\s in {P31/p11,P31/p21,P32/p12,P32/p22,P0/x1,P0/x2,P0/x3}
        \draw[red,dashed] (a\s.south) -- (A\l.north);
    \end{scope}
    \begin{scope}[visible on=<4->]
      \foreach \z in {11,12,13,21,22,23,31,32,33} {
        \draw (CP\z) -- (c\z);
        \draw[szary] (CP\z.south) -- (c0.north);
      }
      \draw (CP0) -- (c0);
    \end{scope}
    \node[lab,visible on=<2->] at (-0.7,2.35) {krok (2):\\ kodowanie};
    \node[lab,visible on=<3->] at (-0.7,1.15) {krok (3):\\ rekurencja};
    \node[lab,visible on=<4->] at (-0.7,0) {krok (4):\\ dekodowanie};
  \end{tikzpicture}
  \vspace{0.3em}
  \begin{itemize}\footnotesize
    \item<2-> Każdy $A_\lambda$ łączy co najwyżej 3 slice'y (linia ciągła: $+$, przerywana: $-$), np.
      $A_{P_{31}}=a_{x_3}-a_{p_{11}}-a_{p_{21}}$, \ $A_{P_0}=-a_{x_1}-a_{x_2}-a_{x_3}$. To samo robimy z $b$ i $\psi_\lambda$.
    \item<3-> Dziesięć par mnożymy wywołaniami rekurencyjnymi $C_\lambda=\Full(A_\lambda,B_\lambda)$.
    \item<4-> $c_{z_{ij}}=C_{P_{ij}}$, a $c_{z_0}$ to suma wszystkich dziesięciu, bo każdy składnik wnosi wkład do $z_0$.
  \end{itemize}
\end{frame}

\begin{frame}{Drzewo rekursji $\Full$}
  \begin{itemize}
    \item Wywołanie osiągnięte przez wybór składników $\tau_1,\dots,\tau_k$ na poziomach $1,\dots,k$ nazywamy
      \textbf{wierzchołkiem} $\tau_1\cdots\tau_k$ (napis $k$ składników). Leży on na głębokości $k$,
      a jego tablice wejściowe są na napisach długości $L-k$.
    \item Korzeń ma $k=0$. \textbf{Liście} $\tau=\tau_1\cdots\tau_L$ (głębokość $L$) to $10^L$ mnożeń z kroku (1).
  \end{itemize}
  \centering
  \begin{tikzpicture}[v/.style={draw,rounded corners=2pt,fill=blue!8,font=\scriptsize,inner sep=2pt},
      lf/.style={draw,circle,fill=orange!40,inner sep=1.2pt}]
    \node[v] (r) at (5.4,0) {korzeń ($k=0$)};
    \foreach \t [count=\i] in {P_{11},P_{12},P_{13},P_{21},P_{22},P_{23},P_{31},P_{32},P_{33},P_0} {
      \node[v] (c\i) at ({(\i-1)*1.2},-1.2) {$\t$};
      \draw (r.south) -- (c\i.north);
    }
    \foreach \t [count=\j] in {P_{11},P_{12},P_{13},P_{21},P_{22},P_{23},P_{31},P_{32},P_{33},P_0}
      \draw[szary] (c3.south) -- ({2.4+(\j-5.5)*0.6},-2.2) node[below,font=\tiny,text=black] {$\t$};
    \node[font=\scriptsize,anchor=west] at (6.2,-2.4) {np.\ wierzchołek $P_{13}P_{21}$ ($k=2$)};
    \node[font=\scriptsize] at (5.4,-3.1) {$\vdots$\quad głębokość $L$: $10^L$ liści $\tau_1\cdots\tau_L$};
  \end{tikzpicture}
\end{frame}

\begin{frame}{Koszt $\Full$}
  \textbf{Koszt.} Na głębokości $k$ jest $10^k$ wierzchołków. Każdy ma dwie tablice wejściowe po $7^{L-k}$
  wpisów i tablicę wyjściową z $10^{L-k}$ wpisami. Krok (2) kosztuje $O(7^{L-k})$ (dziesięć kombinacji co najwyżej
  trzech slice'ów po $7^{L-k-1}$ wpisów), a krok (4) $O(10^{L-k})$:
  \pause
  \[
    \sum_{k=0}^{L}10^k\cdot O\big(7^{L-k}+10^{L-k}\big)=O(L\cdot10^L).
  \]
  \pause
  \begin{itemize}
    \item Kodowanie: $\displaystyle\sum_{k=0}^{L}10^k7^{L-k}=10^L\sum_{k=0}^{L}\big(\tfrac{7}{10}\big)^{L-k}
      <\tfrac{10}{3}\cdot10^L=O(10^L)$.
    \item Dekodowanie: $\displaystyle\sum_{k=0}^{L}10^k10^{L-k}=(L+1)\,10^L$. To ono dominuje.
  \end{itemize}
  \uncover<4->{\idea{$\Full$ kosztuje niewiele więcej niż $10^L$ mnożeń w liściach: kodowanie $O(10^L)$, dekodowanie $O(L\cdot 10^L)$.}}
\end{frame}

\begin{frame}{Co dokładnie liczy $\Full$? Rozwinięcie rekurencji}
  \begin{definicja}
    Dla liścia $\tau$:
    \[
      \Phi_\tau(a):=\sum_u a[u]\prod_{\ell=1}^{L}\varphi_{\tau_\ell}(u_\ell),\qquad
      \Psi_\tau(b):=\sum_v b[v]\prod_{\ell=1}^{L}\psi_{\tau_\ell}(v_\ell),
    \]
    sumy po wszystkich lewych napisach $u$ i prawych napisach $v$ długości $L$.
  \end{definicja}
  \pause
  \begin{itemize}
    \item $\Phi_\tau(a)$ to liczba, którą kodowanie dostarcza do liścia $\tau$: na poziomie $\ell$ stosujemy formę $\varphi_{\tau_\ell}$.
    \item Przykład $L=1$: $\Phi_\lambda(a)=\sum_s\varphi_\lambda(s)\,a[s]$, czyli forma $\varphi_\lambda$ wyliczona na siedmiu liczbach $a[s]$.
    \item Przykład $L=2$: $\Phi_{P_0P_{11}}(a)=\sum_{s,s'}\varphi_{P_0}(s)\varphi_{P_{11}}(s')\,a[ss']
      =-\sum_{i=1}^{3}\big(a[x_ix_1]+a[x_ip_{11}]\big)$.
  \end{itemize}
\end{frame}

\begin{frame}{Wkład wierzchołka do napisu i tablica $\Mult$}
  \begin{definicja}
    Wierzchołek $\tau_1\cdots\tau_k$ \textbf{wnosi wkład} do napisu wyjściowego $w$, jeśli $\tau_\ell$ wnosi wkład
    do $w_\ell$ dla każdego poziomu $\ell\le k$. Zatem liść $\tau$ wnosi wkład do $w$ wtedy i tylko wtedy, gdy $\tau_\ell=P_{ij}$
    na każdym poziomie $\ell$, na którym $w_\ell=z_{ij}$. Tam, gdzie $w_\ell=z_0$, $\tau_\ell$ jest dowolne.
    \[
      \Mult(a,b)[w]:=\sum_{\tau\text{ wnosi wkład do }w}\Phi_\tau(a)\,\Psi_\tau(b).\tag{2}
    \]
  \end{definicja}
  \pause
  \pause
  \begin{przyklad}[$L=2$]
    Do $w=z_{13}z_0$ wnosi wkład 10 liści $P_{13}\lambda$ (dowolne $\lambda$). Do $w=z_{13}z_{21}$ tylko liść $P_{13}P_{21}$.
    Do $w=z_0z_0$ wszystkie $100$ liści.
  \end{przyklad}
\end{frame}

\begin{frame}{Lemat 7}
  \begin{lemat}{7}
    W liściu $\tau$ algorytm $\Full(a,b)$ mnoży $\Phi_\tau(a)$ przez $\Psi_\tau(b)$ i zwraca $\Mult(a,b)$.
  \end{lemat}
  \pause
  \begin{proof}\small
    \setlength{\abovedisplayskip}{3pt}\setlength{\belowdisplayskip}{3pt}
    Indukcja po $L$.
    \textbf{$L=0$:} jedynym liściem jest pusty napis $\tau$. Iloczyny po poziomach są puste, więc
    $\Phi_\tau(a)=a$, $\Psi_\tau(b)=b$, a krok (1) zwraca $ab=\Mult(a,b)$.

    \uncover<3->{\textbf{$L\ge1$, kodowanie:} dla liścia $\lambda\tau'$ (składnik $\lambda$, potem napis $\tau'$ długości $L-1$),
    ponieważ $a_s[u']=a[su']$:
    \[
      \Phi_{\lambda\tau'}(a)=\sum_s\varphi_\lambda(s)\sum_{u'}a[su']\prod_{\ell=1}^{L-1}\varphi_{\tau'_\ell}(u'_\ell)
      =\sum_s\varphi_\lambda(s)\,\Phi_{\tau'}(a_s)=\Phi_{\tau'}\Big(\sum_s\varphi_\lambda(s)a_s\Big)=\Phi_{\tau'}(A_\lambda).
    \]
    Podobnie $\Psi_{\lambda\tau'}(b)=\Psi_{\tau'}(B_\lambda)$.}
    \uncover<4->{Z założenia indukcyjnego liść $\lambda\tau'$ mnoży $\Phi_{\lambda\tau'}(a)$ przez $\Psi_{\lambda\tau'}(b)$ oraz
    $C_\lambda[w']=\sum_{\tau'}\Phi_{\lambda\tau'}(a)\Psi_{\lambda\tau'}(b)$, suma po $\tau'$ wnoszących wkład do $w'$.}

    \uncover<5->{\textbf{Dekodowanie:} do $zw'$ wnoszą wkład dokładnie liście $\lambda\tau'$, gdzie $\lambda$ wnosi wkład do $z$, a $\tau'$ do $w'$.
    Zatem $c[zw']=\sum_{\lambda\text{ wnosi wkład do }z}C_\lambda[w']$ zgadza się z (2).}
  \end{proof}
\end{frame}

\begin{frame}{Współczynniki $\gamma$ i Lemat 8}
  \small
  Dla zmiennej lewej $s$, prawej $t$ i wyjściowej $z$ niech
  \[
    \gamma(s,t,z):=\sum_{\lambda\text{ wnosi wkład do }z}\varphi_\lambda(s)\,\psi_\lambda(t).\tag{3}
  \]
  To współczynnik przy jednomianie $stz$ po lewej stronie tożsamości. Z Lematu 6 jest to też współczynnik przy $stz$ w $G+E$, więc $\gamma(s,t,z)\in\{0,\pm1\}$.
  \pause
  \begin{lemat}{8}
    Dla wszystkich tablic $a,b$ i każdego napisu wyjściowego $w$:
    \[
      \Mult(a,b)[w]=\sum_{u,v}\gamma(u_1,v_1,w_1)\,\gamma(u_2,v_2,w_2)\cdots\gamma(u_L,v_L,w_L)\;a[u]\,b[v],
    \]
    gdzie suma przebiega po wszystkich lewych napisach $u$ i prawych napisach $v$.
  \end{lemat}
\end{frame}

\begin{frame}{Dowód Lematu 8}
  \begin{proof}
    Z definicji $\Phi_\tau$, $\Psi_\tau$:
    \(\Phi_\tau(a)\Psi_\tau(b)=\sum_{u,v}a[u]b[v]\prod_{\ell=1}^{L}\varphi_{\tau_\ell}(u_\ell)\psi_{\tau_\ell}(v_\ell).\)
    \pause
    Zatem współczynnik przy $a[u]b[v]$ w $\Mult(a,b)[w]$ to
    \[
      \sum_{\tau\text{ wnosi wkład do }w}\ \prod_{\ell=1}^{L}\varphi_{\tau_\ell}(u_\ell)\psi_{\tau_\ell}(v_\ell).
    \]
    \pause
    Liście wnoszące wkład do $w$ to te, których składnik na poziomie $\ell$ wnosi wkład do $w_\ell$,
    \textbf{niezależnie dla każdego $\ell$}. Suma iloczynów jest więc iloczynem sum:
    \[
      \prod_{\ell=1}^{L}\ \sum_{\tau_\ell\text{ wnosi wkład do }w_\ell}\varphi_{\tau_\ell}(u_\ell)\psi_{\tau_\ell}(v_\ell)
      \overset{(3)}{=}\prod_{\ell=1}^{L}\gamma(u_\ell,v_\ell,w_\ell).\qedhere
    \]
  \end{proof}
  \pause
  \small\textbf{Dwa opisy $\Mult$:}
  \begin{itemize}
    \item Lemat 8 posłuży do odczytania iloczynów macierzy (2.3.3).
    \item Lemat 7 (suma po liściach) posłuży do pomijania liści, które nie wnoszą wkładu do chcianych wpisów (2.4).
  \end{itemize}
\end{frame}

\begin{frame}{Zbiór wewnętrzny, część wewnętrzna i zewnętrzna}
  \begin{definicja}
    \begin{itemize}
      \item \textbf{Zbiór wewnętrzny} napisu to zbiór poziomów, na których ma on zmienną wewnętrzną:
        $\Wew{p}$ (zamiast $\Zew{x}$) dla lewego, $\Wew{q}$ (zamiast $\Zew{y}$) dla prawego, $\Wew{z_0}$ (zamiast $\Zew{z_{ij}}$) dla wyjściowego.
      \item \textbf{Część wewnętrzna} to napis jego zmiennych na poziomach zbioru wewnętrznego (w kolejności poziomów),
        a \textbf{część zewnętrzna} to napis zmiennych na pozostałych poziomach.
    \end{itemize}
    Napis jest wyznaczony przez swój zbiór wewnętrzny, część zewnętrzną i część wewnętrzną.
  \end{definicja}
  \pause
  \centering
  \begin{tikzpicture}[c/.style={draw,minimum width=0.8cm,minimum height=0.5cm,font=\small}]
    \foreach \v [count=\i] in {x_1,p_{11},x_2,x_3,p_{22},x_1} {
      \ifnum\i=2 \node[c,fill=tlowew] at (\i*0.8,0) {$\v$};
      \else\ifnum\i=5 \node[c,fill=tlowew] at (\i*0.8,0) {$\v$};
      \else \node[c] at (\i*0.8,0) {$\v$};\fi\fi
      \node[font=\tiny,szary] at (\i*0.8,0.5) {\i};
    }
    \node[anchor=east] at (0.3,0) {$u=$};
    \node[anchor=west,font=\small,align=left] at (5.3,0)
      {zbiór wewnętrzny $Q=\{2,5\}$, część wewnętrzna \Wew{$p_{11}p_{22}$},\\ część zewnętrzna \Zew{$x_1x_2x_3x_1$}\quad ($L=6$)};
  \end{tikzpicture}
\end{frame}

\begin{frame}{Parametry i indeksowanie macierzy napisami}
  \small
  Ustalmy $m\ge1$ oraz
  \[
    N_0:=3^{L-m},\qquad D:=4^m.
  \]
  \begin{itemize}
    \item Z $2^L$ iloczynów (po jednym na zbiór $Q\subseteq\{1,\dots,L\}$ poziomów ze skalarnym) używamy tylko
      $\binom Lm$ tych z $|Q|=m$. Wejścia pozostałych ustawiamy na zero.
    \item Iloczyn dla takiego $Q$ to macierz $N_0\times D$ razy $D\times N_0$: $L-m$ poziomów zewnętrznych
      (po 3 zmienne) daje $N_0=3^{L-m}$ wierszy i kolumn, a $m$ poziomów wewnętrznych (po 4 zmienne) daje
      $D=4^m$ składników sumy.
  \end{itemize}
  \pause
  \small
  \textbf{Indeksowanie:}
  \begin{itemize}
    \item Macierz $N_0\times D$: wiersze to napisy z $L-m$ \Zew{zewnętrznych lewych} zmiennych ($x$),
      kolumny to napisy z $m$ \Wew{wewnętrznych lewych} ($p$).
    \item Macierz $D\times N_0$: wiersze to napisy z $m$ \Wew{wewnętrznych prawych} ($q$),
      kolumny to napisy z \Zew{zewnętrznych prawych} ($y$).
    \item Kolumnę $\pi=p_{i_1j_1}\cdots p_{i_mj_m}$ łączymy z wierszem $\pi'=q_{i_1j_1}\cdots q_{i_mj_m}$
      (te same indeksy):
      \[
        (AB)[r,c]:=\sum_\pi A[r,\pi]\,B[\pi',c],
      \]
      suma po wszystkich $D$ napisach $\pi$. Wiersze $r$ wyniku to napisy z $x$, a kolumny $c$ to napisy z $y$.
  \end{itemize}
\end{frame}

\begin{frame}{Kodowanie macierzy w tablicy}
  Niech $Q\subseteq\{1,\dots,L\}$, $|Q|=m$, i $K:=\binom Lm$.
  \begin{itemize}
    \item Lewe napisy $u$ o zbiorze wewnętrznym $Q$ indeksują wpisy macierzy $X_Q$ wymiaru $N_0\times D$:
      $X_Q[u]$ to wpis w wierszu \Zew{(część zewnętrzna $u$)} i kolumnie \Wew{(część wewnętrzna $u$)}.
    \item Prawe napisy $v$ o zbiorze wewnętrznym $Q$ indeksują $Y_Q$ wymiaru $D\times N_0$: wiersz to \Wew{część wewnętrzna}, kolumna to \Zew{część zewnętrzna}.
    \item Napis wyjściowy $w$ o zbiorze wewnętrznym $Q$ ze zmiennymi $z_{i_1j_1},\dots,z_{i_{L-m}j_{L-m}}$ poza $Q$ wskazuje
      wpis $(X_QY_Q)[w]$ w wierszu $x_{i_1}\cdots x_{i_{L-m}}$ i kolumnie $y_{j_1}\cdots y_{j_{L-m}}$.
  \end{itemize}
  \pause
  \begin{block}{Wejście: macierze ułożone obok siebie}
    Dla każdego z $K$ zbiorów $Q$ bierzemy \emph{dowolne} $X_Q$ ($N_0\times D$) i $Y_Q$ ($D\times N_0$) i ustawiamy
    \[
      a[u]:=X_Q[u],\qquad b[v]:=Y_Q[v]
    \]
    dla napisów o zbiorze wewnętrznym $Q$ z $|Q|=m$, a $a[u]:=0$, $b[v]:=0$ dla wszystkich pozostałych napisów.
  \end{block}
\end{frame}

\begin{frame}{Od napisów do wpisów macierzy ($L=6$, $m=2$)}
  \centering
  \begin{tikzpicture}[c/.style={draw,minimum width=0.72cm,minimum height=0.48cm,font=\scriptsize,inner sep=1pt}]
    \foreach \row/\name/\vals in {0/u/{x_1,p_{11},x_2,x_3,p_{22},x_1},
                                   -0.8/v/{y_3,q_{11},y_1,y_3,q_{22},y_2},
                                   -1.6/w/{z_{13},z_0,z_{21},z_{33},z_0,z_{12}}} {
      \node[anchor=east,font=\small] at (0.25,\row) {$\name$:};
      \foreach \v [count=\i] in \vals {
        \ifnum\i=2 \node[c,fill=tlowew] at (\i*0.72,\row) {$\v$};
        \else\ifnum\i=5 \node[c,fill=tlowew] at (\i*0.72,\row) {$\v$};
        \else \node[c] at (\i*0.72,\row) {$\v$};\fi\fi
      }
    }
    \node[font=\scriptsize,anchor=west,align=left] at (4.8,0) {\Zew{$x_1x_2x_3x_1$}: wiersz $u$ w $X_Q$,\quad \Wew{$p_{11}p_{22}$}: kolumna};
    \node[font=\scriptsize,anchor=west,align=left] at (4.8,-0.8) {\Wew{$q_{11}q_{22}$}: wiersz $v$ w $Y_Q$,\quad \Zew{$y_3y_1y_3y_2$}: kolumna};
    \node[font=\scriptsize,anchor=west,align=left] at (4.8,-1.6) {wiersz \Zew{$x_1x_2x_3x_1$} i kolumna \Zew{$y_3y_1y_3y_2$} odczytane z $z_{ij}$};
    \node[font=\scriptsize] at (2.5,-2.2) {poziomy $\ell=1,\dots,6$; cieniowane: $Q=\{2,5\}$};
    % macierze
    \begin{scope}[shift={(1,-6.2)},visible on=<2->]
      \draw (0,0) rectangle (0.8,3); \node[font=\scriptsize] at (0.4,3.3) {$X_Q$: $81\times16$};
      \fill[blue!30] (0,2.1) rectangle (0.8,2.25);
      \node[font=\tiny,anchor=east] at (0,2.17) {wiersz $x_1x_2x_3x_1$};
      \node at (1.3,1.5) {$\cdot$};
      \draw (1.8,2.2) rectangle (4.8,3); \node[font=\scriptsize] at (3.3,3.3) {$Y_Q$: $16\times81$};
      \fill[blue!30] (3.9,2.2) rectangle (4.05,3);
      \node[font=\tiny] at (3.97,1.95) {kolumna $y_3y_1y_3y_2$};
      \node at (5.3,1.5) {$=$};
      \draw (5.8,0) rectangle (8.8,3); \node[font=\scriptsize] at (7.3,3.3) {$X_QY_Q$: $81\times81$};
      \fill[chciane] (6.0+1.9,2.1) rectangle (6.15+1.9,2.25);
      \draw[chciane,-{Stealth}] (8.2,2.17) -- (9.3,2.17) node[right,font=\scriptsize] {$(X_QY_Q)[w]$};
    \end{scope}
  \end{tikzpicture}
  \vspace{0.2em}

  \uncover<3->{Jedno uruchomienie $\Full$ liczy wszystkie $\binom62=15$ iloczynów $X_QY_Q$, po jednym na każdy $Q$ rozmiaru 2.}
\end{frame}

\begin{frame}{Przykład $L=2$, $m=1$: gdzie leżą macierze w tablicach}
  $N_0=3$, $D=4$, $K=\binom21=2$: dwa iloczyny $3\times4$ razy $4\times3$, dla $Q=\{1\}$ i $Q=\{2\}$.
  \begin{columns}[T]
    \begin{column}{0.42\textwidth}
      \centering
      \begin{tikzpicture}[scale=0.36]
        \fill[black!10] (0,4) rectangle (3,7);
        \fill[orange!35] (3,4) rectangle (7,7);
        \fill[green!30] (0,0) rectangle (3,4);
        \fill[black!10] (3,0) rectangle (7,4);
        \draw[step=1,szary] (0,0) grid (7,7);
        \draw[thick] (0,0) rectangle (7,7);
        \foreach \s [count=\i] in {x_1,x_2,x_3,p_{11},p_{12},p_{21},p_{22}} {
          \node[font=\tiny,anchor=east] at (0,{7.5-\i}) {$\s$};
          \node[font=\tiny,rotate=90,anchor=west] at ({\i-0.5},7.1) {$\s$};
        }
        \node[font=\scriptsize] at (1.5,5.5) {$0$};
        \node[font=\scriptsize] at (5,5.5) {$X_{\{2\}}$};
        \node[font=\scriptsize] at (1.5,2) {$X_{\{1\}}^{\!\top}$};
        \node[font=\scriptsize] at (5,2) {$0$};
        \node[font=\scriptsize] at (3.5,-0.9) {tablica $a$: wiersz $u_1$, kolumna $u_2$};
      \end{tikzpicture}
    \end{column}
    \begin{column}{0.56\textwidth}
      \centering
      \begin{tikzpicture}[scale=0.3]
        \fill[black!10] (0,1) rectangle (9,10);
        \fill[orange!35] (9,1) rectangle (10,10);
        \fill[green!30] (0,0) rectangle (9,1);
        \fill[black!10] (9,0) rectangle (10,1);
        \draw[step=1,szary] (0,0) grid (10,10);
        \draw[thick] (0,0) rectangle (10,10);
        \foreach \z [count=\i] in {z_{11},z_{12},z_{13},z_{21},z_{22},z_{23},z_{31},z_{32},z_{33},z_0} {
          \node[font=\tiny,anchor=east] at (0,{10.5-\i}) {$\z$};
          \node[font=\tiny,rotate=90,anchor=west] at ({\i-0.5},10.1) {$\z$};
        }
        \node[font=\tiny,align=center] at (4.5,5.5) {$|Q|\ne m$\\ (nie używamy)};
        \node[font=\scriptsize,anchor=west] at (10.3,5.5) {$X_{\{2\}}Y_{\{2\}}$};
        \node[font=\scriptsize] at (4.5,-0.9) {$X_{\{1\}}Y_{\{1\}}$};
        \node[font=\scriptsize] at (5,-2.4) {tablica wyników: wiersz $w_1$, kolumna $w_2$};
      \end{tikzpicture}
    \end{column}
  \end{columns}
  \vspace{0.2em}
  \small
  \begin{itemize}
    \item $u=x_ip_{jk}$ ma zbiór wewnętrzny $\{2\}$ i wiersz $x_i$, kolumnę $p_{jk}$ w $X_{\{2\}}$. Napis $u=p_{jk}x_i$ ma $\{1\}$, a $x_i$ jest teraz na poziomie 2 (stąd transpozycja na rysunku).
    \item Napis $w=z_{ij}z_0$ to wpis $(X_{\{2\}}Y_{\{2\}})[x_i,y_j]$, a $w=z_0z_{ij}$ to wpis $(X_{\{1\}}Y_{\{1\}})[x_i,y_j]$.
  \end{itemize}
\end{frame}

\begin{frame}{Lemat 9: $\Full$ naprawdę mnoży macierze}
  \begin{lemat}{9}
    Dla tablic wejściowych $a,b$ zdefiniowanych powyżej i każdego napisu wyjściowego $w$, którego zbiór wewnętrzny $Q$
    ma dokładnie $m$ elementów,
    \[
      \Mult(a,b)[w]=(X_QY_Q)[w].\tag{4}
    \]
  \end{lemat}
  \pause
  \begin{proof}[Dowód (1/2)]
    Z Lematu 8: $\Mult(a,b)[w]=\sum_{u,v}\gamma(u_1,v_1,w_1)\cdots\gamma(u_L,v_L,w_L)\,a[u]b[v]$, gdzie
    $\gamma(s,t,z)$ to współczynnik przy $stz$ w $G+E$. Sprawdzamy, które składniki mogą być niezerowe.
    \begin{enumerate}
      \item<3-> $a[u]b[v]\ne0$ wymaga, by zbiory wewnętrzne $u$ i $v$ miały \textbf{dokładnie $m$} elementów.
      \item<4-> \textbf{Poziom $\ell\in Q$:} $w_\ell=z_0$. Jedyne jednomiany $G+E$ z $z_0$ to $p_{ij}q_{ij}z_0$ (współczynnik 1),
        bo $E$ zawiera tylko $z_{ij}$. Zatem $(u_\ell,v_\ell)=(p_{ij},q_{ij})$ dla pewnych $i,j$.
        Czyli $u$ i $v$ mają zmienną wewnętrzną na każdym poziomie $Q$, a ich zbiory wewnętrzne
        (rozmiaru $m$) są \textbf{równe $Q$}.
    \end{enumerate}
  \end{proof}
\end{frame}

\begin{frame}{Dowód Lematu 9 (2/2) i dlaczego $E$ nie szkodzi}
  \begin{proof}[Dowód (2/2)]
    \begin{enumerate}\setcounter{enumi}{2}
      \item \textbf{Poziom $\ell\notin Q$:} $w_\ell=z_{ij}$, a $u_\ell,v_\ell$ są zewnętrzne (zbiory wewnętrzne to dokładnie $Q$).
        Jedyny jednomian $G+E$ z $z_{ij}$ \emph{bez zmiennej wewnętrznej} to $x_iy_jz_{ij}$ (współczynnik 1).
        Zatem $(u_\ell,v_\ell)=(x_i,y_j)$.
      \item<2-> Zostają składniki, w których $u$ ma zbiór wewnętrzny $Q$, część zewnętrzną równą wierszowi $r$ napisu $w$
        i jakąś część wewnętrzną $\pi$, a $v$ ma zbiór wewnętrzny $Q$, część zewnętrzną równą kolumnie $c$ napisu $w$
        i część wewnętrzną $\pi'$ o tych samych indeksach co $\pi$. Jest jeden taki składnik na każdy z $D$ napisów $\pi$,
        ze współczynnikiem 1:
        \[
          \sum_\pi X_Q[r,\pi]\,Y_Q[\pi',c]=(X_QY_Q)[w].\qedhere
        \]
    \end{enumerate}
  \end{proof}
\end{frame}

\begin{frame}{Dlaczego $E$ nie szkodzi: podsumowanie 2.3.3}
  \idea{Każdy jednomian $E$ ma \Wew{wewnętrzną zmienną wejściową} przy \Zew{zewnętrznym $z_{ij}$}. Na poziomie spoza $Q$
    dałoby to $u$ lub $v$ o zbiorze wewnętrznym rozmiaru $>m$, a tam $a[u]b[v]=0$. \textbf{Błąd $E$ ginie.}}
  \vspace{0.8em}
  \begin{przyklad}[$L=2$, $m=1$, $w=z_{21}z_0$, czyli $Q=\{2\}$]
    Jednomian $x_2\qh_{21}z_{21}=x_2q_{21}z_{21}$ z $E$ daje $\gamma(x_2,q_{21},z_{21})=1$. Ale na poziomie 2 ($w_2=z_0$)
    napis $v$ musi mieć $q_{jk}$, więc $v=q_{21}q_{jk}$ ma zbiór wewnętrzny $\{1,2\}$ rozmiaru $2\ne m$.
    Zatem $b[v]=0$ i ten składnik sumy z Lematu 8 znika.
  \end{przyklad}
  \uncover<2->{Podsumowanie: jedno uruchomienie $\Full$ ($10^L$ mnożeń) liczy wszystkie $K$ iloczynów $X_QY_Q$.
    Mogą być zupełnie niezależne. Za chwilę wstawimy tę samą macierz do różnych kopii $X_Q$.}
\end{frame}

\begin{frame}{Kafelkowanie: iloczyn $N\times D\times N$ z iloczynów $N_0\times D\times N_0$}
  \begin{columns}[T]
    \begin{column}{0.55\textwidth}
      \footnotesize
      \begin{itemize}
        \item Tniemy $X$ na \textbf{bloki wierszy} po $N_0$ kolejnych wierszy, a $Y$ na \textbf{bloki kolumn} po $N_0$ kolumn.
          $XY$ składa się z iloczynów blok wierszy $\times$ blok kolumn, każdy kształtu $N_0\times D$ razy $D\times N_0$ (krótko: $N_0\times D\times N_0$).
        \item<2-> $K_0:=\lfloor\sqrt K\rfloor$. Grupujemy bloki w \textbf{pasy} po $K_0$ kolejnych bloków.
          \textbf{Kafelek} to pas wierszy wraz z pasem kolumn. Każdy kafelek to \textbf{jedno uruchomienie $\Full$}.
        \item<3-> Ustalamy $K_0^2\le K$ różnych zbiorów $Q$ rozmiaru $m$ (na rysunku $Q_1,\dots,Q_9$), po jednym na iloczyn blokowy w siatce,
          \emph{takie same w każdym kafelku}. $X_Q$, $Y_Q$ to blok wierszy i blok kolumn iloczynu blokowego z tym $Q$,
          a dla pozostałych $Q$: $X_Q=Y_Q=0$.
        \item<4-> Lemat 9: jedno $\Full$ liczy wszystkie $K_0^2$ iloczynów kafelka. Wszystkie kafelki dają całe $XY$.
      \end{itemize}
    \end{column}
    \begin{column}{0.43\textwidth}
      \centering
      \begin{tikzpicture}[scale=0.5]
        % X
        \draw (0,0) rectangle (0.7,6);
        \foreach \i in {1,...,5} \draw[szary] (0,\i) -- (0.7,\i);
        \draw[very thick] (0,3) -- (0.7,3);
        \node[font=\scriptsize] at (0.35,6.4) {$X$};
        % Y
        \draw (1.2,6.5) rectangle (7.2,7.2);
        \foreach \i in {1,...,5} \draw[szary] ({1.2+\i},6.5) -- ({1.2+\i},7.2);
        \draw[very thick] (4.2,6.5) -- (4.2,7.2);
        \node[font=\scriptsize] at (7.6,6.85) {$Y$};
        % XY
        \draw[step=1,szary,shift={(1.2,0)}] (0,0) grid (6,6);
        \draw[very thick] (1.2,0) rectangle (7.2,6);
        \draw[very thick] (4.2,0) -- (4.2,6); \draw[very thick] (1.2,3) -- (7.2,3);
        \node[font=\scriptsize,below] at (4.2,0) {$N=6N_0$, \ $K_0=3$: 4 kafelki};
        \only<2->{
          \fill[red!15] (0,0) rectangle (0.7,3);
          \fill[red!15] (1.2,6.5) rectangle (4.2,7.2);
          \fill[red!15] (1.2,0) rectangle (4.2,3);
          \draw[step=1,szary,shift={(1.2,0)}] (0,0) grid (3,3);
          \draw[very thick,chciane] (1.2,0) rectangle (4.2,3);
        }
        \only<3->{
          \foreach \r in {0,1,2} \foreach \c in {0,1,2} {
            \pgfmathtruncatemacro{\q}{3*\r+\c+1}
            \node[font=\tiny] at ({1.7+\c},{2.5-\r}) {$Q_{\q}$};
          }
        }
        \node<2->[font=\scriptsize,align=left,anchor=north west] at (0,-0.9)
          {kafelek: pas wierszy 2 $\times$ pas kolumn 1,\\ $K_0^2=9$ iloczynów blokowych z jednej\\
           lewej tablicy (pas wierszy) i jednej prawej (pas kolumn)};
      \end{tikzpicture}
    \end{column}
  \end{columns}
\end{frame}

\begin{frame}{Dlaczego siatka i ile to kosztuje}
  \begin{itemize}
    \item Bloki parujemy w siatkę, a nie dowolnie, więc \textbf{lewa tablica kafelka zależy tylko od jego pasa wierszy},
      a prawa tylko od pasa kolumn. Pozwoli to \alert{współdzielić kodowania} między kafelkami (2.4.1).
    \item Najpierw dopełniamy $N$ do wielokrotności $K_0N_0$ zerowymi wierszami $X$ i kolumnami $Y$.
      To co najwyżej podwaja $N$, bo $N\ge K_0N_0$ (pokażemy w 2.4.4).
  \end{itemize}
  \pause
  \textbf{Koszt.} Niech $M:=KN_0^2$ będzie liczbą wpisów wyniku we wszystkich $K$ iloczynach kafelka.
  Ponieważ $K_0\ge\sqrt K/2$:
  \[
    \#\text{kafelków}=\Big(\frac{N}{K_0N_0}\Big)^2=\frac{N^2}{K_0^2N_0^2}\le\frac{N^2}{(K/4)N_0^2}=\frac{4N^2}{M}.
  \]
  \pause
  Razem $O\big(\frac{N^2}{M}\cdot L\cdot10^L\big)$ operacji, czyli $O(L\cdot10^L/M)$ na wpis $XY$.
  \begin{itemize}
    \item Przy $L=10m$ daje to $N^2\operatorname{poly}(L)$, a przy wybranym przez nas $L=19m$ około $N^2D^{0.2}$.
    \item<4-> Mimo to wybieramy $L=19m$. Dlaczego, zobaczymy w 2.4.3 (rzadkość).
    \item<4-> Kafelkowanie powiększa tylko wymiary zewnętrzne, więc obsługuje mniejsze $D$ niż algorytm Coppersmitha z tej samej tożsamości
      ($D\approx N^{0.1204}$ zamiast $N^{0.1402}$ [Cop82]). Za to \alert{każdy wpis $XY$ to jedno wyjście jednego uruchomienia $\Full$}.
  \end{itemize}
\end{frame}

% ======================================================================
\section{2.4 Liczymy tylko chciane wpisy}
% ======================================================================

\begin{frame}{Cel i dwie modyfikacje}
  Chcemy policzyć tylko wpisy $XY$ na pozycjach z $W$, gdzie $|W|\le N^2/\sqrt D$.
  Algorytm z 2.3 modyfikujemy na dwa sposoby:
  \begin{enumerate}
    \item<2-> \textbf{Współdzielenie kodowania (2.4.1).} Kodowanie (krok (2)) każdej tablicy wejściowej zależy
      \emph{tylko od niej}, a każda tablica jest wspólna dla wielu kafelków. Liczymy je \textbf{raz} na pas.
    \item<3-> \textbf{Przycinanie drzewa (2.4.2).} W krokach (3)--(4) pomijamy wywołanie rekurencyjne, jeśli żadne
      z jego wyjść nie prowadzi do wpisu z $W$.
  \end{enumerate}
  \uncover<4->{Potem:
  \begin{itemize}
    \item liczymy odwiedzane liście (2.4.3, \textbf{tu jest nowa analiza}),
    \item sumujemy koszty (2.4.4).
  \end{itemize}}
\end{frame}

\begin{frame}{2.4.1 Kodowanie liczymy raz}
  \begin{itemize}
    \item Liczby $\Phi_\tau(a)$ i $\Psi_\tau(b)$ mnożone w liściu $\tau$ zależą odpowiednio \textbf{tylko od $a$} i \textbf{tylko od $b$}.
      Liczymy je więc raz na tablicę wejściową, a nie raz na kafelek.
    \item<2-> \textbf{Kodowanie} tablicy $a$ to tablica $10^L$ liczb $\Phi_\tau(a)$ indeksowana liśćmi $\tau$.
    \item<3-> Liczymy ją, uruchamiając $\Full$ bez $b$ i bez kroku (4). Z Lematu 7 lewa liczba w liściu $\tau$ to $\Phi_\tau(a)$,
      więc liść ją zapisuje. Koszt to $O(10^L)$, czyli koszt kodowania z 2.3.1.
    \item<3-> Tak samo liczymy kodowanie $b$, czyli liczby $\Psi_\tau(b)$.
    \item<4-> Liczymy kodowania tablic wejściowych \textbf{wszystkich pasów wierszy i kolumn}: $2N/(K_0N_0)$ tablic po $10^L$ liczb.
      Każdy kafelek czyta swoje dwa kodowania.
  \end{itemize}
\end{frame}

\begin{frame}{Dlaczego współdzielenie kodowania jest konieczne}
  \begin{columns}[T]
    \begin{column}{0.56\textwidth}
      \begin{itemize}
        \item Kodowanie ma
          \[10^L=(9+1)^L>\binom Lm9^{L-m}=M\]
          liczb (to jeden wyraz rozwinięcia dwumianu), czyli \emph{więcej niż wpisów kafelka}.
        \item<2-> Kodowanie każdego kafelka osobno kosztowałoby więc więcej niż jedną operację na wpis całego $XY$,
          i to zanim w ogóle skorzystamy z $W$.
        \item<3-> Przy współdzieleniu $2N/(K_0N_0)$ kodowań obsługuje wszystkie $(N/(K_0N_0))^2$ kafelków.
          Każde kosztuje $O(10^L)$, niezależnie od $W$.
        \item<4-> Założenie $N\ge D^{18}$ czyni ten koszt pomijalnym (policzymy w 2.4.4).
      \end{itemize}
    \end{column}
    \begin{column}{0.42\textwidth}
      \centering
      \begin{tikzpicture}[scale=0.55]
        \draw[step=1,szary,shift={(1.2,0)}] (0,0) grid (4,4);
        \draw[very thick] (1.2,0) rectangle (5.2,4);
        \foreach \i in {1,...,4} {
          \node[draw,fill=blue!15,font=\tiny,minimum width=0.6cm] (r\i) at (-0.2,{4.5-\i}) {$\Phi(a^{(\i)})$};
          \node[draw,fill=orange!25,font=\tiny,minimum height=0.6cm,rotate=90] (c\i) at ({0.7+\i},5.4) {$\Psi(b^{(\i)})$};
        }
        \only<3->{
          \fill[blue!12] (1.2,2) rectangle (5.2,3);
          \fill[orange!20] (3.2,0) rectangle (4.2,4);
          \fill[red!25] (3.2,2) rectangle (4.2,3);
          \draw[step=1,szary,shift={(1.2,0)}] (0,0) grid (4,4);
          \draw[very thick] (1.2,0) rectangle (5.2,4);
          \draw[-{Stealth},thick,blue!70] (r2.east) -- (1.25,2.5);
          \draw[-{Stealth},thick,orange!80!black] (3.7,4.35) -- (3.7,3.95);
          \node[font=\tiny,red!70!black] at (3.7,2.5) {$(2,3)$};
        }
        \node[font=\scriptsize,align=center] at (2.4,-1.0) {$a^{(i)}$: lewa tablica pasa wierszy $i$,\\
          $b^{(j)}$: prawa tablica pasa kolumn $j$.\\ Kafelek $(i,j)$ czyta $\Phi(a^{(i)})$ i $\Psi(b^{(j)})$.};
      \end{tikzpicture}
    \end{column}
  \end{columns}
\end{frame}

\begin{frame}{2.4.2 Pomijamy niepotrzebne wywołania}
  \begin{itemize}
    \item Po policzeniu kodowań pomijamy tę fazę $\Full$. Każdy liść czyta swoje dwie liczby z kodowań,
      a każde wywołanie wykonuje już tylko dekodowanie z kroku (4).
    \item Każdemu wywołaniu przekazujemy zbiór $S$ \textbf{wyjść, których od niego chcemy}, czyli prowadzących do wpisów z $W$.
    \item Wyjścia wierzchołka $\tau_1\cdots\tau_k$ są indeksowane napisami wyjściowymi długości $L-k$, więc $S$ jest zbiorem takich napisów.
  \end{itemize}
  \pause
  \begin{definicja}
    Dla zbioru $S$ napisów wyjściowych i zmiennej wyjściowej $z$ \textbf{slice} to
    \[S_z:=\{w' : zw'\in S\}.\]
    Tablica na $S$ przypisuje liczbę całkowitą każdemu napisowi z $S$.
  \end{definicja}
  \pause
  Jeśli $U$ to zbiór przekazany korzeniowi (w końcowym algorytmie są to napisy chcianych wpisów kafelka), to w Lemacie 10 zobaczymy, że $S$
  składa się z sufiksów $w_{k+1}\cdots w_L$ tych napisów $w\in U$, do których wierzchołek wnosi wkład.
\end{frame}

\begin{frame}{Algorytm $\Pruned$}
  \begin{block}{$\Pruned_{\tau_1\cdots\tau_k}(S)$ dla zbioru $S$ napisów wyjściowych długości $L-k$; zwraca tablicę wierzchołka $\tau_1\cdots\tau_k$ algorytmu $\Full$ ograniczoną do $S$}
    \begin{enumerate}
      \item Jeśli $k=L$: odczytaj $\Phi_\tau(a)$ i $\Psi_\tau(b)$ z policzonych kodowań i zwróć $\Phi_\tau(a)\Psi_\tau(b)$ dla liścia $\tau=\tau_1\cdots\tau_L$.
      \item Dla każdego składnika $\lambda$ niech $S_\lambda$ będzie sumą slice'ów $S_z$ po $z$, do których $\lambda$ wnosi wkład:
        \[S_{P_{ij}}:=S_{z_{ij}}\cup S_{z_0},\qquad S_{P_0}:=S_{z_0}.\]
      \item Dla każdego $\lambda$ z $S_\lambda\neq\emptyset$ policz $C_\lambda:=\Pruned_{\tau_1\cdots\tau_k\lambda}(S_\lambda)$, tablicę na $S_\lambda$.
      \item Zwróć tablicę $c$ na $S$ o slice'ach $c_{z_{ij}}:=C_{P_{ij}}$ i $c_{z_0}:=\sum_\lambda C_\lambda$,
        ograniczonych do $S_{z_{ij}}$ i $S_{z_0}$.
    \end{enumerate}
  \end{block}
  \pause
  \begin{itemize}
    \item Ograniczenia w (4) mają sens, bo $S_{z_{ij}}\subseteq S_{P_{ij}}$ i $S_{z_0}\subseteq S_\lambda$ dla każdego $\lambda$ (z kroku (2)).
    \item Różnica względem $\Full$: brak kodowania (jest gotowe) i wywołania tylko dla $S_\lambda\ne\emptyset$.
  \end{itemize}
\end{frame}

\begin{frame}{Animacja: $\Pruned$ dla $L=2$, $U=\{z_{13}z_0,\ z_0z_{21}\}$}
  \centering
  \begin{tikzpicture}[
      v/.style={draw,rounded corners=2pt,font=\scriptsize,inner sep=2pt,minimum width=0.95cm,fill=blue!8},
      s/.style={font=\tiny,text=chciane},
      lf/.style={draw,circle,inner sep=1.3pt,fill=chciane}]
    \node[v,fill=red!15] (r) at (6.3,0) {korzeń};
    \node[s,above=1pt of r] {$S=\{z_{13}z_0,\ z_0z_{21}\}$};
    \node[font=\scriptsize,align=left,anchor=west,visible on=<2->] at (9.0,0.1)
      {$S_{z_{13}}=\{z_0\}$, \ $S_{z_0}=\{z_{21}\}$,\\ pozostałe slice'y puste};
    \foreach \l/\t [count=\i] in {P11/P_{11},P12/P_{12},P13/P_{13},P21/P_{21},P22/P_{22},P23/P_{23},P31/P_{31},P32/P_{32},P33/P_{33},P0/P_0} {
      \node[v,visible on=<3->] (\l) at ({(\i-1)*1.4},-1.35) {$\t$};
      \draw[visible on=<3->] (r.south) -- (\l.north);
    }
    \node[s,below=1pt of P13,visible on=<3->] {$\{z_0,z_{21}\}$};
    \foreach \l in {P11,P12,P21,P22,P23,P31,P32,P33,P0}
      \node[s,below=1pt of \l,visible on=<3->] {$\{z_{21}\}$};
    % poziom 2
    \foreach \l in {P11,P12,P21,P22,P23,P31,P32,P33,P0} {
      \node[lf,visible on=<4->] (L\l) at ($(\l)+(0,-1.35)$) {};
      \draw[visible on=<4->] ($(\l.south)+(0,-0.3)$) -- (L\l);
      \node[font=\tiny,below=1pt of L\l,visible on=<4->] {$P_{21}$};
    }
    \foreach \k in {1,...,10} {
      \node[lf,visible on=<5->] (Q\k) at ({2.8+(\k-5.5)*0.2},-3.0) {};
      \draw[visible on=<5->] ($(P13.south)+(0,-0.3)$) -- (Q\k);
    }
    \node[font=\tiny,align=center,visible on=<5->] at (2.8,-3.4) {wszystkie 10};
  \end{tikzpicture}

  \begin{itemize}\footnotesize
    \item<3-> Krok (2): $S_{P_{13}}=S_{z_{13}}\cup S_{z_0}=\{z_0,z_{21}\}$, a dla $\lambda\ne P_{13}$: $S_\lambda=S_{z_0}=\{z_{21}\}$.
      Wywołujemy wszystkie 10 dzieci.
    \item<4-> Pod $\lambda\ne P_{13}$ jest $S=\{z_{21}\}$: niepuste jest tylko $S_{P_{21}}=\{\pusty\}$ ($\pusty$ to pusty napis). Jeden liść $\lambda P_{21}$.
    \item<5-> Pod $P_{13}$ jest $S=\{z_0,z_{21}\}$: $S_{z_0}=\{\pusty\}$, więc wszystkie 10 liści $P_{13}\lambda$.
    \item<6-> \alert{Odwiedzonych liści: $10+9=19$ ze $100$.} Rzeczywiście: do $z_{13}z_0$ wnosi wkład 10 liści, do $z_0z_{21}$
      też 10, a wspólny jest jeden ($P_{13}P_{21}$). Łączny rozmiar zbiorów: $2+11+19=32\le(L+1)\cdot19$.
  \end{itemize}
\end{frame}

\begin{frame}{Lemat 10: $\Pruned$ jest poprawny i odwiedza tylko potrzebne liście}
  Przypomnienie (2): $\Mult(a,b)[w]$ to suma $\Phi_\tau(a)\Psi_\tau(b)$ po liściach wnoszących wkład do $w$, czyli wybierających
  $P_{ij}$ na każdym poziomie, gdzie $w$ ma $z_{ij}$, i dowolny składnik tam, gdzie $w$ ma $z_0$.
  \begin{definicja}
    Dla zbioru $U$ napisów wyjściowych długości $L$, $\Leaves(U)$ to zbiór liści wnoszących wkład do \emph{jakiegoś} napisu z $U$.
  \end{definicja}
  \pause
  \begin{lemat}{10}
    Niech $a$ i $b$ będą tablicami wejściowymi, a $U$ zbiorem napisów wyjściowych długości $L$. Wywołany w korzeniu,
    $\Pruned(U)$ zwraca $\Mult(a,b)$ ograniczone do $U$. Odwiedza dokładnie liście z $\Leaves(U)$, a zbiory
    przekazywane do jego wywołań mają łączny rozmiar co najwyżej $(L+1)\,|\Leaves(U)|$.
  \end{lemat}
\end{frame}

\begin{frame}{Dowód Lematu 10}
  \begin{proof}\small
    \textbf{Wartości:} indukcja po $L-k$. Każde $C_\lambda$ to tablica odpowiedniego dziecka w $\Full$ ograniczona do $S_\lambda$,
    więc krok (4) to krok (4) algorytmu $\Full$ ograniczony do $S$.

    \uncover<2->{\textbf{Odwiedzane liście:} $S_\lambda$ składa się z ostatnich $L-k-1$ zmiennych tych napisów z $S$, do których
    \emph{pierwszej} zmiennej $\lambda$ wnosi wkład. Indukcja po $k$ pokazuje, że zbiór przekazany wierzchołkowi $\tau_1\cdots\tau_k$
    to sufiksy $w_{k+1}\cdots w_L$ napisów $w\in U$, do których ten wierzchołek wnosi wkład, i że wierzchołek jest wywołany
    \emph{wtedy i tylko wtedy}, gdy ten zbiór jest niepusty. Zatem liść jest odwiedzony $\iff$ wnosi wkład do jakiegoś napisu z $U$.}

    \uncover<3->{\textbf{Łączny rozmiar:} każdemu sufiksowi w zbiorze wierzchołka $\tau_1\cdots\tau_k$ przypisujemy liść,
    który przedłuża $\tau_1\cdots\tau_k$, wybierając $P_{ij}$ tam, gdzie sufiks ma $z_{ij}$, i $P_0$ tam, gdzie ma $z_0$.
    Wierzchołek wnosi wkład do pewnego $w\in U$ z tym sufiksem, a wtedy ten liść także, więc liść należy do $\Leaves(U)$.}
    \uncover<4->{Liść wyznacza wierzchołek (prefiks długości $k$) i sufiks (odczytany z jego składników). Zbiory na każdej
    z $L+1$ głębokości mają więc łączny rozmiar co najwyżej $|\Leaves(U)|$.}
  \end{proof}
  \uncover<5->{\idea{Każde wywołanie jest wykonane \textbf{co najwyżej raz}, niezależnie od tego, ile napisów z $U$ z niego korzysta.
    Czas zależy od \textbf{rozmiaru sumy} zbiorów liści, a nie od sumy ich rozmiarów. Trzeba oszacować tę sumę.}}
\end{frame}

% ----------------------------------------------------------------------
\subsection{2.4.3 Liczenie liści}

\begin{frame}{2.4.3 Ile liści wnosi wkład do jednego napisu?}
  Od teraz $U$ zawiera tylko napisy wyjściowe długości $L$ o zbiorze wewnętrznym rozmiaru dokładnie $m$
  (wskazują wpisy $X_QY_Q$, Lemat 9). Jest ich
  \[M=\binom Lm9^{L-m}=KN_0^2.\]
  \pause
  \begin{itemize}
    \item Do jednego napisu o zbiorze wewnętrznym $Q$ wnosi wkład $10^m=D^{\log_410}\approx D^{1.66}$ liści: dowolny składnik na każdym z $m$ poziomów $Q$
      i $P_{ij}$ na każdym poziomie poza $Q$.
    \item<3-> To \alert{więcej niż $D$} mnożeń potrzebnych do policzenia wpisu wprost jako iloczynu skalarnego!
    \item<4-> Dlatego \textbf{nie} wolno liczyć liści każdego napisu osobno. Większość liści napisu wnosi wkład
      także do wielu innych napisów, a takich liści jest łącznie niewiele.
  \end{itemize}
\end{frame}

\begin{frame}{Prywatny liść i rząd liścia}
  \begin{definicja}
    \textbf{Prywatny liść} napisu wyjściowego to ten z jego liści, który wybiera $P_0$ na każdym poziomie $Q$.
  \end{definicja}
  Nie wnosi on wkładu do żadnego innego napisu. Napis o innym zbiorze wewnętrznym wymaga składnika $\ne P_0$ na jakimś
  poziomie $Q$. Napis o tym samym $Q$, ale innej części zewnętrznej (inne $z_{ij}$ na jakimś poziomie poza $Q$), wymaga tam innego składnika.
  \pause
  \begin{definicja}
    \textbf{Rząd} liścia to $m$ minus liczba poziomów, na których wybiera on $P_0$.
  \end{definicja}
  \begin{itemize}
    \item Liść wnoszący wkład do napisu wybiera $P_0$ tylko na poziomach zbioru wewnętrznego napisu, więc jego rząd jest $\ge0$.
      Rząd to liczba poziomów $Q$, na których liść wybiera któryś z $P_{ij}$, współdzielonych z iloczynem zewnętrznym.
    \item Jedynym liściem rzędu $0$ wnoszącym wkład do napisu jest jego prywatny liść.
  \end{itemize}
\end{frame}

\begin{frame}{Liczby $\alpha_d$ i $\beta_d$}
  \[
    \alpha_d:=\binom md9^d\qquad\text{liści rzędu $d$ wnosi wkład do \emph{jednego} napisu}
  \]
  (prywatny liść, w którym na $d$ poziomach $Q$ zamieniamy $P_0$ na jeden z 9 składników $P_{ij}$).
  \pause
  \[
    \beta_d:=\binom L{m-d}9^{L-m+d}\qquad\text{liści rzędu $d$ jest w \emph{całym} drzewie}
  \]
  (wybieramy $m-d$ poziomów z $P_0$ i jeden z 9 pozostałych składników na każdym innym poziomie).
  W szczególności $\beta_0=\binom Lm9^{L-m}=M$, bo liście rzędu 0 to dokładnie prywatne liście, po jednym na napis.
\end{frame}

\begin{frame}{Podwójne zliczanie}
  \begin{itemize}
    \item Liść rzędu $d$ wnosi wkład do $\binom{L-m+d}{d}$ napisów, po jednym na każdy zbiór wewnętrzny rozmiaru $m$
      zawierający $m-d$ poziomów, na których liść wybiera $P_0$. (Dobieramy $d$ z $L-(m-d)$ pozostałych poziomów; część zewnętrzna
      wynika z $P_{ij}$ liścia.)
    \item<2-> Liczymy pary (napis, liść rzędu $d$ wnoszący do niego wkład) na dwa sposoby:
      \[
        M\cdot\alpha_d=\beta_d\cdot\binom{L-m+d}{d}.
      \]
  \end{itemize}
\end{frame}

\begin{frame}{Przykład: liście jednego napisu ($L=6$, $m=2$)}
  Napis $w=z_{13}\,\Wew{z_0}\,z_{21}\,z_{33}\,\Wew{z_0}\,z_{12}$, \ $Q=\{2,5\}$. Liść to napis składników wybranych na kolejnych poziomach.
  \vspace{0.4em}
  \centering
  \begin{tikzpicture}[c/.style={draw,minimum width=0.85cm,minimum height=0.5cm,font=\scriptsize}]
    \foreach \row/\vals/\lab in {
      0/{P_{13},P_0,P_{21},P_{33},P_0,P_{12}}/{rząd 0: prywatny liść, $\alpha_0=1$, wnosi wkład do 1 napisu},
      -0.85/{P_{13},P_{ij},P_{21},P_{33},P_0,P_{12}}/{},
      -1.45/{P_{13},P_0,P_{21},P_{33},P_{ij},P_{12}}/{rząd 1: $\alpha_1=\binom21\cdot9=18$ liści, każdy wnosi wkład do $\binom51=5$ napisów},
      -2.3/{P_{13},P_{ij},P_{21},P_{33},P_{ij},P_{12}}/{rząd 2: $\alpha_2=\binom22\cdot81=81$ liści, każdy wnosi wkład do $\binom62=15$ napisów}} {
      \foreach \v [count=\i] in \vals {
        \ifnum\i=2 \node[c,fill=tlowew] at (\i*0.85,\row) {$\v$};
        \else\ifnum\i=5 \node[c,fill=tlowew] at (\i*0.85,\row) {$\v$};
        \else \node[c] at (\i*0.85,\row) {$\v$};\fi\fi
      }
      \node[anchor=west,font=\scriptsize] at (5.6,\row) {\lab};
    }
  \end{tikzpicture}
  \vspace{0.4em}
  \begin{itemize}
    \item Liście wyższego rzędu: \textbf{więcej na napis} ($\alpha_d$ rośnie), ale każdy \textbf{współdzielony przez więcej napisów}.
    \item<2-> Przy $L=19m$ jest ich \textbf{mniej w sumie} (nierówność (5) na następnym slajdzie).
  \end{itemize}
\end{frame}

\begin{frame}{Nierówność (5): liści wysokiego rzędu jest geometrycznie mniej}
  \[
    \frac{\beta_d}{\beta_{d-1}}
    =\frac{\binom{L}{m-d}}{\binom{L}{m-d+1}}\cdot9
    =\frac{9(m-d+1)}{L-m+d}.
  \]
  \pause
  Dla $L=19m$ i $1\le d\le m$: licznik $\le9m$, mianownik $=18m+d\ge18m+1$, więc
  \[
    \frac{\beta_d}{\beta_{d-1}}=\frac{9(m-d+1)}{L-m+d}\le\frac{9m}{18m+1}<\frac12,
    \qquad\text{zatem}\qquad \beta_d\le2^{-d}M.\tag{5}
  \]
  \pause
  \begin{itemize}
    \item Liczymy \textbf{wszystkie} liście rzędu $d>m/9$, niezależnie od tego, czy wnoszą wkład do napisu z $U$.
      To tanie, bo $\beta_d$ maleją geometrycznie.
    \item<4-> To wymaga $m$ \textbf{dużo poniżej} $L/10$. Przy $m\approx L/10$ (wybór, przy którym całe $XY$ liczy się tanio, 2.3.4)
      iloraz byłby bliski 1 dla małych $d$, więc $\beta_d\approx M$ i nie byłoby oszczędności.
    \item<5-> Małe $m$ sprawia też, że $D=4^m$ jest małe względem $N_0=3^{18m}$ (około $N_0^{0.07}$). Kafelek musi mieścić się w iloczynie,
      więc $N$ musi być dużą potęgą $D$. Razem z kosztem kodowań to jest powód założenia \alert{$N\ge D^{18}$}.
  \end{itemize}
\end{frame}

\begin{frame}{Lemat 11: mało liści wnosi wkład do rzadkiego zbioru wpisów}
  Od teraz ustalamy $L=19m$, więc
  \[
    N_0=3^{18m},\qquad K=\binom{19m}{m},\qquad D=4^m,\qquad 2^m=\sqrt D,\qquad 2^{-m/9}=D^{-1/18}.
  \]
  \begin{lemat}{11}
    Dla każdego zbioru $U$ napisów wyjściowych, których zbiory wewnętrzne mają dokładnie $m$ elementów,
    \[
      |\Leaves(U)|\le\sum_{d=0}^{m}\min\{|U|\alpha_d,\ \beta_d\},
    \]
    a jeśli $L=19m$, to
    \[
      |\Leaves(U)|\le2^{-m/9}\big(2^m|U|+2M\big).
    \]
  \end{lemat}
\end{frame}

\begin{frame}{Dowód Lematu 11: pierwsza nierówność}
  \begin{proof}[Dowód (1/2)]
    Liczbę liści rzędu $d$ w $\Leaves(U)$ szacujemy na dwa sposoby:
    \begin{itemize}
      \item przez $|U|\alpha_d$: każdy napis z $U$ „płaci” za swoje liście rzędu $d$;
      \item przez $\beta_d$: liczbę \emph{wszystkich} liści rzędu $d$.
    \end{itemize}
    Każdy liść z $\Leaves(U)$ ma rząd $0\le d\le m$, więc sumując mniejsze z dwóch oszacowań po $d$, dostajemy pierwszą nierówność.
  \end{proof}
  \pause
  \textbf{Interpretacja} (każde z oszacowań osobno):
  \begin{itemize}
    \item $\sum_d|U|\alpha_d=|U|\sum_d\binom md9^d=|U|\cdot10^m$, czyli koszt liczenia każdego napisu z $U$ osobno;
    \item $\sum_d\beta_d<2M$ z (5), czyli koszt policzenia wszystkich $M$ napisów.
    \item Ani $|U|\alpha_d$, ani $\beta_d$ nie zależą od tego, \emph{które} napisy są w $U$.
  \end{itemize}
\end{frame}

\begin{frame}{Dowód Lematu 11: druga nierówność ($L=19m$)}
  \begin{proof}[Dowód (2/2)]
    Dla rzędów $d\le m/9$ używamy $|U|\alpha_d$, a dla $d>m/9$ używamy $\beta_d$.

    \uncover<2->{\textbf{Rzędy $d>m/9$.} Niech $d_0$ będzie najmniejszą liczbą całkowitą $>m/9$. Z (5):
    \[
      \sum_{d>m/9}\beta_d\le M\sum_{d\ge d_0}2^{-d}=2\cdot2^{-d_0}M<2\cdot2^{-m/9}M.
    \]}
    \uncover<3->{\textbf{Rzędy $d\le m/9$.} Dają co najwyżej $|U|\sum_{d\le m/9}\binom md9^d$. Ponieważ
    $9^d=72^d\,8^{-d}\le72^{m/9}\,8^{-d}$ dla $d\le m/9$,
    \[
      \sum_{d\le m/9}\binom md9^d\le72^{m/9}\sum_{d=0}^{m}\binom md8^{-d}
      =72^{m/9}\Big(\frac98\Big)^m=\Big(72\cdot\Big(\frac98\Big)^9\Big)^{m/9}<256^{m/9}=2^m\cdot2^{-m/9},
    \]
    bo $72\cdot(9/8)^9\approx207{,}83<208<256$.}

    \uncover<4->{Dodając obie części: $|\Leaves(U)|\le2^{-m/9}\big(2^m|U|+2M\big)$.}
  \end{proof}
\end{frame}

\begin{frame}{Rysunek 7: oszacowanie z Lematu 11 (schematycznie)}
  \centering
  \begin{tikzpicture}[x=0.9cm,y=0.9cm]
    % schemat jak Rysunek 7 w pracy: m/9 w x=4, m w x=8; pionowo skala logarytmiczna
    \fill[blue!15,visible on=<2->] (0,0) -- (0,1.3) -- (4,2.7) -- (4,0) -- cycle;
    \fill[red!15,visible on=<3->] (4,0) -- (4,2.15) -- (8,0.6) -- (8,0) -- cycle;
    \draw[->] (0,0) -- (8.6,0) node[right,font=\small] {rząd $d$};
    \draw[->] (0,0) -- (0,4.4) node[above,font=\small] {liczba liści (skala log.)};
    \draw[blue!80!black,thick] (0,1.3) -- (8,4.1) node[right,font=\small] {$|U|\,\alpha_d$};
    \draw[red!80!black,thick] (0,3.7) -- (8,0.6) node[right,font=\small] {$\beta_d\le2^{-d}M$};
    \node[left,font=\small] at (0,3.7) {$M$};
    \node[left,font=\small] at (0,1.3) {$|U|$};
    \draw[dashed,visible on=<2->] (4,0) -- (4,2.7);
    \node[below,font=\small] at (4,0) {$m/9$};
    \node[below,font=\small] at (8,0) {$m$};
    \node[align=center,font=\scriptsize,visible on=<2->] at (2,0.75) {małe rzędy:\\ płaci każdy\\ napis z $U$};
    \node[align=center,font=\scriptsize,visible on=<3->] at (6,0.55) {duże rzędy:\\ bierzemy wszystkie liście};
  \end{tikzpicture}

  \vspace{0.4em}
  \begin{itemize}\small
    \item Liści rzędu $d$ w $\Leaves(U)$ jest co najwyżej $\min\{|U|\alpha_d,\beta_d\}$.
      $|U|\alpha_d$ rośnie z $d$ aż do $d\approx0{,}9m$ (do każdego napisu wnosi wkład więcej liści wyższego rzędu),
      a $\beta_d$ maleje geometrycznie (liści wyższego rzędu jest mniej).
    \item<2-> Dowód dzieli sumę w $d=m/9$, co nie musi być punktem przecięcia obu oszacowań.
    \item<3-> Całość szacujemy przez zacieniowane pole.
  \end{itemize}
\end{frame}

\begin{frame}{Co mówi Lemat 11 dla gęstości z Twierdzenia 5}
  Ponieważ $2^m=\sqrt D$ i $2^{-m/9}=D^{-1/18}$, druga nierówność to
  \[
    |\Leaves(U)|\le D^{-1/18}\big(\sqrt D\,|U|+2M\big).
  \]
  \pause
  Weźmy kafelek, w którym chcemy $|U|=M/\sqrt D$ wpisów (gęstość z Twierdzenia 5: $N^2/\sqrt D$ z $N^2$). Wtedy $\sqrt D|U|=M$ i
  \[
    |\Leaves(U)|\le D^{-1/18}\cdot3M=O\big(M/D^{1/18}\big).
  \]
  \pause
  \begin{itemize}
    \item To o czynnik $D^{1/18}$ \textbf{mniej niż $M$} wpisów kafelka.
    \item I dużo mniej niż $|U|\cdot D=M\sqrt D$ mnożeń przy liczeniu każdego chcianego wpisu wprost.
    \item<4-> Oszacowanie jest \textbf{liniowe w $|U|$}, więc po zsumowaniu po kafelkach wynik nie zależy od tego, jak chciane wpisy
      rozkładają się na kafelki.
  \end{itemize}
\end{frame}

% ======================================================================
\section{2.4.4 Dowód Twierdzenia 5}
% ======================================================================

\begin{frame}{Algorytm}
  \begin{enumerate}
    \item Policz kodowania wszystkich pasów wierszy i pasów kolumn (2.4.1).
    \item<2-> Każda chciana pozycja $(I,J)\in W$ leży w \textbf{jednym} kafelku $T$ i jest wskazana przez \textbf{jeden} napis
      wyjściowy kafelka $T$. Jego zbiór wewnętrzny to zbiór $Q$ iloczynu blokowego zawierającego $(I,J)$, a wiersz i kolumna to wiersz
      i kolumna $(I,J)$ wewnątrz tego iloczynu blokowego (2.3.3, 2.3.4).
    \item<3-> Dla kafelka $T$ niech $W_T$ będzie zbiorem napisów wyjściowych chcianych pozycji w $T$.
    \item<4-> Dla każdego kafelka $T$ z $W_T\ne\emptyset$ uruchom $\Pruned(W_T)$.
    \item<5-> Dla każdego $(I,J)\in W$ zwróć wartość przy jego napisie wyjściowym.
  \end{enumerate}
  \uncover<6->{\vspace{0.5em}\textbf{Poprawność:} z Lematu 9 (napis wskazuje właściwy wpis iloczynu blokowego, czyli $(XY)[I,J]$) i Lematu 10
    ($\Pruned$ zwraca $\Mult$ na $W_T$).}
  \uncover<7->{
  \begin{center}
    \begin{tikzpicture}[b/.style={draw,rounded corners,fill=blue!8,font=\scriptsize,align=center,minimum height=0.9cm},
        arr/.style={-{Stealth},thick}]
      \node[b] (e) {kodowania\\ pasów};
      \node[b,right=0.6cm of e] (w) {$W\to W_T$\\ (napisy)};
      \node[b,right=0.6cm of w,fill=red!12] (p) {$\Pruned(W_T)$\\ dla każdego $T$};
      \node[b,right=0.6cm of p] (o) {odczyt\\ $(XY)[I,J]$};
      \draw[arr] (e)--(w); \draw[arr] (w)--(p); \draw[arr] (p)--(o);
    \end{tikzpicture}
  \end{center}}
\end{frame}

\begin{frame}{Konsekwencja (a) założenia $N\ge D^{18}$}
  Ustalamy $D=4^m$ i $L=19m$. Mamy $D^{18}=4^{18m}$, więc potęgi $m$-te porównujemy, porównując ich podstawy.
  \pause
  \begin{block}{(a) Kafelkowanie: $K_0N_0\le N$, więc dopełnienie co najwyżej podwaja $N$}
    Z ogólnego oszacowania $\binom Lm\le(eL/m)^m$: $K=\binom{19m}{m}\le(19e)^m$. Zatem
    \[
      KN_0\le\big(19e\cdot3^{18}\big)^m\le\big(4^{18}\big)^m,\qquad\text{bo}\quad
      19e\cdot3^{18}\approx2{,}001\cdot10^{10}<4^{18}\approx6{,}872\cdot10^{10},
    \]
    a stąd $K_0N_0\le KN_0\le D^{18}\le N$. Dalej $N$ oznacza dopełniony rozmiar.
  \end{block}
\end{frame}

\begin{frame}{Konsekwencja (b) założenia $N\ge D^{18}$}
  \begin{block}{(b) Kodowania są tanie}
    Wykonujemy $2N/(K_0N_0)$ kodowań po $O(10^L)$ operacji. Żeby zmieścić się w budżecie $O(2^{-m/9}N^2)$, potrzeba
    \[
      10^L\le2^{-m/9}\,N\,N_0.\tag{6}
    \]
    Zachodzi:
    $\dfrac{10^L\,2^{m/9}}{N_0}=\Big(\dfrac{10^{19}\cdot2^{1/9}}{3^{18}}\Big)^m\le(4^{18})^m=D^{18}\le N$, \
    bo $\dfrac{10^{19}\cdot2^{1/9}}{3^{18}}\approx2{,}788\cdot10^{10}<4^{18}$.
  \end{block}
  \pause
  Wtedy łączny koszt kodowań to
  \[
    \frac{2N}{K_0N_0}\cdot O(10^L)\overset{(6)}{\le}\frac{2N}{K_0N_0}\cdot O\big(2^{-m/9}NN_0\big)=O\Big(\frac{2^{-m/9}N^2}{K_0}\Big)=O\big(2^{-m/9}N^2\big).
  \]
\end{frame}

\begin{frame}{Koszt}
  \begin{itemize}
    \item \textbf{Kodowania:} z (6)
      \[
        \frac{2N}{K_0N_0}\cdot O(10^L)\le\frac{2N}{K_0N_0}\cdot O\big(2^{-m/9}NN_0\big)=O\big(2^{-m/9}N^2\big).
      \]
    \item<2-> \textbf{Chciane wpisy:} każda pozycja z $W$ trafia do jednego $W_T$, więc
      $\sum_T|W_T|\le|W|\le N^2/\sqrt D=2^{-m}N^2$. Kafelków jest co najwyżej $4N^2/M$, więc $\sum_TM\le4N^2$.
    \item<3-> \textbf{Jedno wywołanie $\Pruned$ ze zbiorem $S$:} $O(L|S|)$ operacji. Każdy krok obsługuje każdy napis z $S$ stałą liczbę razy,
      a napis ma długość $L$. Zbiory trzymamy jako posortowane listy, więc slice to spójny fragment, a suma slice'ów to scalanie.
    \item<4-> Z Lematu 10 (łączny rozmiar zbiorów $\le(L+1)|\Leaves|$) wychodzi $O(L^2)=O(\log^2D)$ operacji \textbf{na liść} z $\Leaves(W_T)$.
  \end{itemize}
  \uncover<5->{
  Z Lematu 11 zsumowanego po kafelkach:
  \[
    O\Big(L^2\sum_T|\Leaves(W_T)|\Big)
    \le O\Big(L^2\,2^{-m/9}\Big(\underbrace{2^m|W|}_{\le N^2}+\underbrace{2\textstyle\sum_TM}_{\le8N^2}\Big)\Big)
    =O\big(L^2\,2^{-m/9}N^2\big),
  \]
  \alert{niezależnie od tego, jak $W$ rozkłada się na kafelki}.}
\end{frame}

\begin{frame}{Koszt: drobiazgi i wynik końcowy}
  \begin{itemize}
    \item \textbf{Drobiazgi:} wypisanie $K_0^2$ zbiorów $Q$, znalezienie napisów wyjściowych i sortowanie zbiorów $W_T$ kosztują $O(L)$
      na zbiór $Q$, na kafelek i na pozycję z $W$. Razem to $O(L\,2^{-m/9}N^2)$, bo:
      \begin{itemize}
        \item $K\le N$,
        \item kafelków jest co najwyżej $4N^2/N_0$,
        \item $|W|\le2^{-m}N^2$,
        \item $N,N_0\ge2^{m/9}$.
      \end{itemize}
  \end{itemize}
  \pause
  \begin{block}{Razem}
    Ponieważ $L=19\log_4D$ i $2^{-m/9}=D^{-1/18}$:
    \[
      O\big(2^{-m/9}N^2\big)+O\big(L^2\,2^{-m/9}N^2\big)+O\big(L\,2^{-m/9}N^2\big)
      =O\!\left(\frac{N^2\log^2D}{D^{1/18}}\right).\qquad\blacksquare
    \]
  \end{block}
\end{frame}

\begin{frame}{Rozmiar słowa}
  Każda liczba w algorytmie jest liczbą całkowitą o module $N^{O(1)}$:
  \begin{itemize}
    \item Wpisy wejścia mają taki rozmiar z założenia.
    \item<2-> Każda liczba w kodowaniu (lub policzona po drodze do niego) to kombinacja $\pm1$ co najwyżej $7^L$ wpisów wejścia,
      bo krok (2) tylko dodaje i odejmuje slice'y.
    \item<3-> Każda wartość policzona przez $\Pruned$ to suma co najwyżej $10^L$ iloczynów dwóch zakodowanych liczb (Lemat 7).
    \item<4-> Z (6): $7^L\le10^L\le2^{-m/9}NN_0\le N^2$.
    \item<5-> Indeksy i pozycje na listach nie przekraczają liczby używanych słów, czyli też $N^{O(1)}$.
  \end{itemize}
  \uncover<6->{Zatem każda liczba ma $O(\log N)$ bitów, czyli mieści się w $O(1)$ słowach modelu word RAM.}
\end{frame}

\begin{frame}{Uwaga 12: związek z innymi technikami}
  \begin{uwagaP}{12}
    \begin{itemize}
      \item Prace o mnożeniu macierzy zwykle liczą tylko \emph{mnożenia}. Tutaj dodawania trzeba kontrolować starannie.
      \item Kodowania (2.4.1) to algorytm Yatesa dla potęg Kroneckera [Yat37, Goo58], standardowy w rekurencyjnych algorytmach
        mnożenia macierzy.
      \item Przycięta rekursja (2.4.2) to analog przycinania FFT [Mar71, SB93], przyciętej inwersji Möbiusa [BHKK10] oraz
        (po transpozycji) rzadkiej ewaluacji obwodów potęg Kroneckera [AL25], opartej na [Wil24].
      \item \textbf{Nowa} w tym kierunku jest analiza z 2.4.3: przycięta rekursja odwiedza \textbf{mało liści}.
    \end{itemize}
  \end{uwagaP}
\end{frame}

% ======================================================================
\section{Podsumowanie}
% ======================================================================

\begin{frame}{Kluczowe idee}
  \begin{enumerate}
    \item<1-> \textbf{Rekurencja à la Strassen} z tożsamością Schönhagego (10 mnożeń zamiast 13):
      $L$ poziomów daje naraz $\binom Lm$ cienkich iloczynów $N_0\times D\times N_0$.
    \item<2-> \textbf{Błąd $E$ ginie}, bo jego jednomiany mieszają zmienne wewnętrzne z zewnętrznymi (Lemat 9).
    \item<3-> \textbf{Kafelkowanie siatką}: każdy wpis $XY$ to jedno wyjście jednego $\Full$, a tablica wejściowa zależy tylko od pasa.
    \item<4-> \textbf{Kodowanie raz na pas}: kodowanie zależy tylko od $X$ albo tylko od $Y$, a $N\ge D^{18}$ czyni je pomijalnym.
    \item<5-> \textbf{Przycinanie}: odwiedzamy tylko $\Leaves(U)$. Każde wywołanie raz, $O(L^2)$ na liść (Lemat 10).
    \item<6-> \textbf{Liczenie liści według rzędu}: prywatne liście ($d=0$) plus liście współdzielone.
      Małe rzędy liczymy na napis, a duże globalnie ($\beta_d$ maleje geometrycznie przy $L=19m$).
      Stąd $O(M/D^{1/18})$ liści na kafelek przy gęstości $1/\sqrt D$.
  \end{enumerate}
\end{frame}

\begin{frame}[c]
  \centering
  \definecolor{claude}{HTML}{D97757}
  \definecolor{diabel}{HTML}{8E44AD}
  \begin{tikzpicture}[scale=1.25]
    % rogi jak w emoji z diabłem (rysowane pod promieniami)
    \foreach \s in {1,-1}
      \fill[diabel,xscale=\s] (0.16,0.70) .. controls (0.45,0.92) and (0.72,1.12) .. (0.84,1.50)
        .. controls (1.02,1.18) and (0.96,0.80) .. (0.60,0.50) -- cycle;
    \foreach \kat/\dl in {74/1.00,116/1.02,146/0.94,180/0.95,212/0.90,239/0.92,
                          272/0.95,311/1.00,343/0.88,5/0.88,36/0.96,95/0.80}
      \fill[claude] ($(0,0)+(\kat+90:0.035)$) -- ($(\kat:\dl)+(\kat+90:0.075)$)
        arc[start angle=\kat+90,end angle=\kat-90,radius=0.075]
        -- ($(0,0)+(\kat-90:0.035)$) -- cycle;
    \fill[claude] (0,0) circle (0.12);
  \end{tikzpicture}

  \vspace{0.8em}
  {\Huge\textcolor{structure}{\textbf{Koniec}}}

  \vspace{1em}
  {\large Stanisław Macura \quad Michał Chłopicki}
\end{frame}


\end{document}
